\documentclass{../../Template/TemplateLecture}

\renewcommand{\LectureName}{\textsc{Einführung in die \sout{Geometrie und} Topologie}}
\renewcommand{\ProfName}{\textsc{Prof. Dr. Markus Hausmann}}
\renewcommand{\Semester}{Sommer 2024}
\renewcommand{\mName}{Daniel Emse}

\renewcommand{\listfigurename}{Liste der Definitionen}

\begin{document}

\pagenumbering{roman}
% \section*{Abkürzungen und Konventionen}
% \subsection*{Abkürzungen}
% \begin{acronym}
%     \acro{VR}[VR]{Vektorraum}
%     \acro{KVR}[$\IK$-VR]{$\IK$-Vektorraum}
%     \acro{KE}[KE]{Körpererweiterung}
%     \acro{GM}[GM]{Gramsche Matrix}
%     \acro{UR}[UR]{Unterraum}
%     \acro{lu}[lu]{linear unabhängig}
%     \acro{ES}[ES]{Erzeugendensystem}
%     \acro{BLF}[BLF]{Bilinearform}
%     \acro{Iso}[Iso]{Isomorphismus}
%     \acro{Endo}[Endo]{Endomorphismus}
%     \acro{lin}[lin]{linear}
%     \acro{bilin}[bilin]{bilinear}
%     \acro{klin}[$k$-lin]{$k$-linear}
%     \acro{EW}[EW]{Eigenwert}
%     \acro{EV}[EV]{Eigenvektor}
%     \acro{ER}[ER]{Eigenraum}
%     \acro{ONB}[ONB]{Orthonormalbasis}
% \end{acronym}
% \subsection*{Konventionen}\addcontentsline{toc}{section}{Abkürzungen und Konventionen}
% \begin{itemize}
%     \item $0 \in \IN$
%     \item $T^0(V) = \IK$
% \end{itemize}

\listoftodos
\subsubsection*{Anmerkungen und Korrekturen}
Falls du Anmerkungen/Korrekturen hast, melde dich bei mir (\href{mailto:daniel.emse@uni-bonn.de}{\texttt{daniel.emse@uni-bonn.de}}). 
\pagebreak
\tableofcontents
\pagebreak
\pagenumbering{arabic}
\setcounter{chapter}{-1}

%%%%%%%%%%%%%%%%%%%%%%%%%%%%%%%%%%%%%%%%%%%%%%%%%%%%%%%%%%%%%%%%%%%%%%%%%%%%%%%%%%%%%%%%%%%%%%%%%%%%
%%%%%%%%%%%%%%%%%%%%%%%%%%%%%%%%%%%%%%%%%% 2024 / 04 / 09 %%%%%%%%%%%%%%%%%%%%%%%%%%%%%%%%%%%%%%%%%%
%%%%%%%%%%%%%%%%%%%%%%%%%%%%%%%%%%%%%%%%%%%%%%%%%%%%%%%%%%%%%%%%%%%%%%%%%%%%%%%%%%%%%%%%%%%%%%%%%%%%

\part{Mengentheoretische Topologie}
\chapter{Einführung}

\breakline{Vorlesung 01 fiel aus; Ersatz durch Videos. Bombenstimmung!}
\emph{Topologie} beschäftigt sich mit dem Studium von Räumen und Formen, zum Beispiel ein \emph{Kreis}, ein \emph{Torus} oder ein \emph{Sierpi\'nski-Raum}.

\subsection*{Grundbegriffe}
\begin{itemize}
    \item topologischer Raum, stetige Abbildung
    \item zusammenhängend
    \item kompakt
    \item Verklebekonstruktion
    \item Fundamentalgruppe (algebraische Invariante)
\end{itemize}

\section*{Motivation}
\begin{defi}{Topologische Äquivalenz, homöomorph}{topEquiv}
    Zwei Teilmengen $X \subseteq \IR^n, Y \subseteq \IR^m$ heißen \emph{topologisch äquivalent}\index{topologische Äquivalenz} (auch \emph{homöomorph}\index{Homöomorphie}), wenn es stetige Abbildungen $f: X \to Y$ und $g: Y \to X$ gibt, sodass $g \circ f = \id_X$ und $f \circ g = \id_Y$.
\end{defi}

\begin{beispiel}\leavevmode
    \begin{enumerate}[$(i)$]
        \item Das Intervall $\intervaloo{-1,1}$ und $\IR$ sind topologisch äquivalent.
        
        \emph{Wähle $f: \IR \to \intervaloo{-1,1}, x \mapsto \frac{x}{\sqrt{1 + x^2}}$ und $g: \intervaloo{-1,1} \to \IR, y \mapsto \frac{y}{\sqrt{1 - y^2}}$.}
        \item Kreis $\IS^1 = \set{x \in \IR^2 \mid \sqrt{x_1^2 + x_2^2} = 1}$ und Quadrat $Q = \set{y \in \IR^2 \mid \max\set{\abs{y_1}, \abs{y_2}} = 1}$ sind topologisch äquivalent.
        
        \emph{Projiziere die Punkte aus dem Kreis durch den Strahl, der sie mit dem Ursprung verbindet, auf das Quadrat und andersrum}
        \item Ein Donut und eine Kaffeetasse sind topologisch äquivalent.\todo{wie geht das}
    \end{enumerate}
\end{beispiel}

\subsection*{Wie zeigt man, dass zwei Räume $X$ und $Y$ \emph{nicht} topologisch äquivalent sind?}
\begin{beispiel}
    $\IR^2$ und $\IR^3$ sind nicht topologisch äquivalent, oder allgemeiner:
\end{beispiel}

\begin{thm}{Satz der topologischen Invarianz der Dimension}{dimesnsionInvariance}
    Für $m \neq n$ sind $\IR^m$ und $\IR^n$ nicht topologisch äquivalent.
\end{thm}

\textbf{Vorsicht:} Es gibt stetige, surjektive Abbildungen $\IR \to \IR^n$ für alle $n$ (\enquote{Peano-Kurven}).

\begin{thm}{$\IR^n$ und $\IR$ sind nicht topologisch äquivalent}{RRnNotEquiv}
    Für $n \geq 2$ sind $\IR$ und $\IR^n$ nicht topologisch äquivalent.
\end{thm}

\emph{Idee:} $\IR \setminus \set{\pt}$ hat zwei zusammenhängende Teile, aber $\IR^n \setminus \set{\pt}$ ist für $n \geq 2$ zusammenhängend.

\begin{defi}{Weg, wege-äquivalent}{path}
    Sei $X \subseteq \IR^n$. Ein \emph{Weg}\index{Weg} in $X$ ist eine stetige Abbildung $w: \intervalcc{0,1} \to X$. Zwei Punkte $x, y \in X$ heißen \emph{wege-äquivalent}\index{wege-äquivalent}, falls es einen Weg in $X$ gibt mit $w(0) = x$ und $w(1) = y$.
\end{defi}

\begin{lem}{Äquivalenzrelation Wege-Äquivalenz}{pathEquivalentER}
    Wege-Äquivalenz ist eine Äquivalenzrelation auf $X$.
\end{lem}

\emph{Beweisidee:} Für Reflexivität betrachte den konstanten Weg, für Symmetrie den Weg in die umgekehrte Richtung und für Transitivität das Hintereinanderlaufen zweier Wege mit doppelter Geschwindigkeit.

\begin{defi}{$\pi_0(X)$}{piZero}
    Für $X \subset \IR^n$ ist $\pi_0(X)$ die \emph{Menge der Weg-Zusammenhangskomponenten}\index{$\pi_0(X)$ (Menge Weg-ZHK)} von Punkten in $X$.
\end{defi}

\begin{beispiel}\leavevmode
    \begin{enumerate}[$(i)$]
        \item $\pi_0(\IR^n)$ hat 1 Element für $n \geq 1$: $x, y \in \IR^n$ sind verbunden durch den linearen Weg $\intervalcc{0,1} \ni t \mapsto x + t \cdot (y - x) \in \IR^n$.
        \item Für $n \geq 2$ hat $\pi_0(\IR^n \setminus \set{p})$ 1 Element.
        \begin{proof}
            \emph{Fall 1:} Die Gerade durch $x$ und $y$ enthält $p$ nicht. Dann liegt der lineare Weg ganz in $\IR^n \setminus \set{p}$, verbindet also $x$ und $y$ in $\IR^n \setminus \set{p}$.

            \emph{Fall 2:} $p$ liegt auf der Geraden von $x$ nach $y$. Dann wähle ein $z \in \IR^n$, dass nicht auf der Geraden durch $x$ und $y$ liegt. Nach Fall 1 ist $x$ wege-äquivalent zu $z$, und $z$ ist wege-äquiavlent zu $y$. Nach Transitivität und Lemma~\ref{lem:pathEquivalentER} ist also $x$ wege-äquiavlent zu $y$.
        \end{proof}
        \item $\pi_0(\IR \setminus \set{0})$ hat genau 2 Elemente.
        \begin{proof}
            Wenn $x, y \in \IR \setminus \set{0}$ beide positiv oder beide negativ sind, so werden sie durch den linearen Weg in $\IR \setminus \set{0}$ verbunden. Also gibt es höchstens zwei Äquivalenzklassen.

            Nach dem Mittelwertsatz verläuft jeder stetige Weg in $\IR$ von einem negativen zu eniem positiven Wert durch 0. Folglich gibt ses keinen Weg in $\IR \setminus \set{0}$ von einer negativen zu einer positiven Zahl. Also hat $\pi_0(\IR \setminus \set{0})$ mindestens zwei Elemente.
        \end{proof}
    \end{enumerate}
\end{beispiel}

\begin{proof}[Beweis von Satz~$\ref{thm:RRnNotEquiv}$]
    Angenommen, $f: \IR \to \IR^n$ und $g: \IR^n \to \IR$ sind stetig und invers zueinander. Setze $p := f(0) \in \IR^n$, dann sind auch $f: \IR \setminus \set{0} \to \IR^n \setminus \set{f(0)}$ und $g: \IR^n \setminus \set{f(0)} \to \IR \setminus \set{0}$ stetig und invers zueinander.

    Nach Beispiel $(ii)$ hat $\IR^n \setminus \set{p}$ nur eine Wegkomponente, es gibt also einen Weg $w$ in $\IR^n \setminus \set{p}$ von $f(-1)$ zu $f(1)$. Betrachte jetzt den zusammengesetzten Weg

    \[
        \begin{tikzcd}
            \intervalcc{0,1} \arrow[r, "w"] \arrow[rr, bend right, "w' := g \circ w"] & \IR^n \setminus \set{p} \arrow[r, "g"] & \IR \setminus \set{0}.
        \end{tikzcd}
    \]
    
    Dann ist $w'$ ein Weg von $w'(0) = g(\overset{f(-1)}{\overbrace{w(0)}}) = -1$ zu $w'(1) = g(\overset{f(1)}{\overbrace{w(1)}}) = 1. \quad \lightning$ Dies ist ein Widerspruch zu Beispiel $(iii)$.
\end{proof}

Um $\IR^2$ und $\IR^n$ mit $n \geq 3$ zu unterscheiden, benutzt man ein höherdimensionales Analogon dieser Idee: \emph{Deformationsklassen von Schleifen} (statt Wegklassen von Punkten).

Sei $\IS := \set{x \in \IR^2 \mid \norm{x} = 1}$.

\begin{defi}{Schlinge, Deformation}{schlinge}
    Eine \emph{Schlinge}\index{Schlinge} in $X$ ist eine stetige Abbildung $f: \IS^1 \to X$. Eine \emph{Deformation}\index{Deformation} von einer Schlinge $f_0$ zu einer Schlinge $f_1$ ist eine stetige Abbildung $D: \IS^1 \times \intervalcc{0,1} \to X$ mit $f_0(x) = D(x, 0), f_1(x) = D(x, 1)$.
\end{defi}

\begin{lem}{Deformation von Schlingen ist Äquivalenzrelation}{deformationSchlinge}
    Wir bezeichnen mit $S(X)$ die Menge der Deformationsklassen von Schlingen in $X$. Später werden wir folgendes zeigen:
    \begin{enumerate}[$(i)$]
        \item Wenn $X$ und $Y$ topologisch äquivalent sind, haben $S(X)$ und $S(Y)$ gleich viele Elemente.
        \item $S(\IR^n \setminus \set{p})$ hat für $n \geq 3$ genau ein Element.
        \item $S(\IR^2 \setminus \set{p})$ hat abzählbar unendlich viele Elemente.
    \end{enumerate}
\end{lem}
Wir werden einen Isomorphismus $S(\IR^2 \setminus \set{0}) \iso \IZ$ konstruieren.


\chapter{Metrische und topologische Räume}
\section{Metrische Räume und offene Mengen}

\setcounter{lemma}{-3}

\textbf{Erinnerung.} Eine \emph{Metrik}\index{Metrik} auf einer Menge $X$ ist eine Funktion $d: X \times X \to \IR_{\geq 0}$ mit den folgenden Eigenschaften:

\begin{enumerate}[$(i)$]
    \item $d(x,y) = 0 \Leftrightarrow x = y$,
    \item $d(x,y) = d(y,x)$ für alle $x,y \in X$,
    \item Dreiecksungleichung $d(x,z) \leq d(x,y) + d(y,z)$ für alle $x,y,z \in X$.
\end{enumerate}
Ein \emph{metrischer Raum}\index{Metrik!-scher Raum} ist ein Paar $(X,d)$ einer Menge $X$ und einer Metrik $d$ darauf.

Seien $(X,d)$ und $(Y,d')$ metrische Räume. Eine Funktion $f: X \to Y$ ist \emph{stetig} in $x \in X$, falls für jedes $\varepsilon > 0$ ein $\delta > 0$ existiert, sodass für alle $x'$ mit $d(x,x') < \delta$ gilt $d(f(x),f(x')) < \varepsilon$.

Die Funktion $f$ heißt \emph{stetig}, wenn sie in jedem Punkt $x \in X$ stetig ist.

\begin{beispiel} \leavevmode
    \begin{enumerate}[$(i)$]
        \item Sei $V$ ein reeller Vektorraum mit Norm $\norm{\cdot}$. Dann definiert $d(v,w) := \norm{v - w}$ eine Metrik auf $V$. Insbesondere: $\IR^n$ mit der euklidischen Norm.
        \item Sei $(X,d)$ ein metrischer Raum und $Y \subseteq X$ eine Teilmenge. Dann ist $(Y,d|_{Y \times Y})$ wieder ein metrischer Raum. Insbesondere ist jedes $Y \subseteq \IR^n$ mit entsprechender Norm ein metrischer Raum.
        \item Sei $X$ eine Menge. Dann ist 
        \begin{equation*}
            d(x,y) = \begin{cases}
                0 & x = y \\
                1 & \text{sonst}
            \end{cases}
        \end{equation*}
        eine Metrik auf $X$.
    \end{enumerate}
\end{beispiel}

Wir wollen den Stetigkeitsbegriff umformen, ohne dabei $d$ explizit zu verwenden.

\textbf{Notation.} Sei $U(x, \varepsilon) := \set{y \in X \mid d(x,y) < \varepsilon}$ der \emph{offene $\varepsilon$-Ball um $x$}.

\textbf{Beobachtung.} Sei $f: X \to Y$ eine Abbildung zwischen metrischen Räumen, $x \in X$ und $\varepsilon, \delta > 0$. Dann sind äquivalent:
\begin{enumerate}[$(i)$]
    \item Für alle $x' \in X$ mit $d(x,x') < \delta$ gilt $d(f(x),f(x')) < \varepsilon$,
    \item Es ist $f(U(x,\delta)) \subseteq U(f(x),\varepsilon)$.
\end{enumerate}

\begin{defi}{Umgebung}{surrounding}
    Sei $(X,d)$ ein metrischer Raum, $U \subseteq X, x \in X$. Dann heißt $U$ \emph{Umgebung von $x$}\index{Umgebung}, falls ein $\varepsilon > 0$ existiert, sodass $U(x, \varepsilon) \subseteq U$.
\end{defi}

\begin{thm}{Stetigkeit $\Leftrightarrow$ Urbilder von Umgebungen sind Umgebung}{stetigOffen}
    Sei $f: X \to Y$ eine Abbildung zwischen metrischen Räumen und $x \in X$. Dann ist $f$ genau dann stetig in $x$, wenn für alle Umgebungen $V$ von $f(x)$ in $Y$ das Urbild $f^{-1}(V)$ eine Umgebung von $x$ ist.
\end{thm}
\begin{proof}
    \begin{itemize}
        \item[\enquote{$\Rightarrow$}] Sei $V$ eine Umgebung von $f(x)$. Dann gibt es ein $\varepsilon > 0$, sodass $U(f(x),\varepsilon) \subseteq V$. Da $f$ stetig in $x$ ist, gibt es ein $\delta > 0$, sodass $f(U(x,\delta)) \subseteq U(f(x),\varepsilon) \subseteq V$. Also ist $U(x,\delta) \subseteq f^{-1}(V)$ und somit ist $f^{-1}(V)$ eine Umgebung von $x$.
        \item[\enquote{$\Leftarrow$}] Sei $\varepsilon > 0$. Dann ist $U(f(x),\varepsilon)$ eine Umgebung von $f(x)$. Also ist $f^{-1}(U(f(x),\varepsilon))$ eine Umgebung von $x$. Somit existiert ein $\delta > 0$ mit $U(x,\delta) \subset f^{-1}(U(f(x),\varepsilon))$. Also ist $f(U(x,\delta)) \subseteq U(f(x),\varepsilon)$, und $f$ ist stetig in $x$.
    \end{itemize}
\end{proof}

\begin{defi}{Offene Menge}{openSet}
    Sei $(X,d)$ ein metrischer Raum. Eine Teilmenge $U \subseteq X$ heißt \emph{offen}\index{offene Menge}, falls sie Umgebung all ihrer Punkte ist.

    \emph{Äquivalent:} Für alle $x \in U$ existiert ein $\varepsilon > 0$ mit $U(x, \varepsilon) \subseteq V$.
\end{defi}

\begin{bemerkung}
    Die offenen $\varepsilon$-Bälle $U(x, \varepsilon)$ sind tatsächlich offen:

    Für $y \in U(x, \varepsilon)$ gilt 
    \begin{equation*}
        U(y, \varepsilon - d(x,y)) \subseteq U(x, \varepsilon)
    \end{equation*}
    aufgrund der Dreiecksungleichung.
\end{bemerkung}

\begin{thm}{Stetigkeit $\Leftrightarrow$ Urbilder offener Mengen sind offen}{continuousOpenSet}
    Sei $f: X \to Y$ eine Abbildung zwischen metrischen Räumen. Dann ist $f$ genau dann stetig, wenn für alle offenen $V \subseteq Y$ das Urbild $f^{-1}(V)$ offen in $X$ ist.
\end{thm}
Man sagt häufig, \emph{Urbilder offener Mengen sind offen}.

\begin{proof}
    \begin{itemize}
        \item[\enquote{$\Rightarrow$}] Sei $V \subseteq Y$ offen und $x \in f^{-1}(V)$. Dann ist $f(x) \in V$, also ist $V$ eine Umgebung von $f(x)$. Nach Satz~\ref{thm:stetigOffen} ist also $f^{-1}(V)$ eine Umgebung von $x$. Also ist $f^{-1}(V)$ offen.
        \item[\enquote{$\Leftarrow$}] Sei $x \in X$ und $V$ eine Umgebung von $f(x)$. Dann enthält $V$ einen $\varepsilon$-Ball $U(f(x),\varepsilon)$ für ein $\varepsilon > 0$. Der Ball $U(x, \varepsilon)$ ist offen in $Y$, also ist $f^{-1}(U(f(x),\varepsilon))$ offen in $X$. Also existiert ein $\delta > 0$, sodass $U(x, \delta) \subseteq f^{-1}(U(f(x),\varepsilon)) \subseteq f^{-1}(V)$. Somit ist $f^{-1}(V)$ eine Umgebung von $x$ und $f$ ist stetig.
    \end{itemize}
\end{proof}


%%%%%%%%%%%%%%%%%%%%%%%%%%%%%%%%%%%%%%%%%%%%%%%%%%%%%%%%%%%%%%%%%%%%%%%%%%%%%%%%%%%%%%%%%%%%%%%%%%%%
%%%%%%%%%%%%%%%%%%%%%%%%%%%%%%%%%%%%%%%%%% 2024 / 04 / 11 %%%%%%%%%%%%%%%%%%%%%%%%%%%%%%%%%%%%%%%%%%
%%%%%%%%%%%%%%%%%%%%%%%%%%%%%%%%%%%%%%%%%%%%%%%%%%%%%%%%%%%%%%%%%%%%%%%%%%%%%%%%%%%%%%%%%%%%%%%%%%%%

\breakline{Beginn Vorlesung 02 (11.04.2024)}

\begin{corollary}
    Stetigkeit hängt nur von den offenen Mengen ab.
\end{corollary}

\begin{lem}{Eigenschaften metrischer Räume}{metricSpaceTraits}
    Sei $(X,d)$ ein metrischer Raum. Dann gilt:
    \begin{enumerate}[$(i)$]
        \item Die leere Menge $\emptyset$ und die Menge $X$ sind offen in $X$.
        \item Sind $U_1, \dots, U_n \subseteq X$ offen, so ist auch $\bigcap_{i = 1}^n U_i$ offen in $X$.
        \item Für jede Familie $\set{U_i}_{i \in I}$ offener Mengen ist auch $\bigcup_{i \in I} U_i$ offen in $X$.
    \end{enumerate}
\end{lem}

\begin{proof}
    \begin{enumerate}[$(i)$]
        \item Folgt aus der Definition offener Mengen.
        \item Sei $x \in \bigcap_{i = 1}^n U_i$. Da jedes $U_i$ offen ist, gibt es ein $\varepsilon_i > 0$, sodass $U(x, \varepsilon_i) \subseteq U_i$. Sei $\varepsilon := \min\set{\varepsilon_1, \dots, \varepsilon_n}$. Dann ist $U(x, \varepsilon) \subseteq U(x, \varepsilon_i) \subseteq U_i$ für alle $i$ und damit gilt $U(x, \varepsilon) \subseteq \bigcap_{i = 1}^n U_i$, diese Menge ist also auch offen.
        \item Sei $x \in \bigcup_{i \in I} U_i$. Dann existiert ein $j \in I$, sodass $x \in U_j$. Da $U_j$ offen ist, existiert ein $\varepsilon > 0$, sodass $U(x, \varepsilon) \subseteq U_j \subseteq \bigcup_{i \in I} U_i$. Damit ist die Vereinigung offen.
    \end{enumerate}
\end{proof}

\begin{bemerkung}
    Die Endlichkeit in $(ii)$ ist wichtig. \emph{Gegenbeispiel:} 
    \begin{equation*}
        \bigcap_{n \in \IN_{>0}} \intervaloo{-\frac{1}{n}, \frac{1}{n}} = \set{0}
    \end{equation*}
    ist nicht offen in $\IR$ mit der Standardmetrik, aber jede Menge $\intervaloo{-\frac{1}{n},\frac{1}{n}}$ ist offen.
\end{bemerkung}

\section{Topologische Räume}
\begin{defi}{Topologie, topologischer Raum}{topology}
    Eine \emph{Topologie}\index{Topologie} auf einer Menge $X$ ist eine Menge $\cO$ von Teilmengen von $X$, die \enquote{offenen Mengen}, sodass gilt:
    \begin{enumerate}[$(i)$]
        \item $\emptyset \in \cO$ und $X \in \cO$.
        \item Sind $U_1, \ldots, U_n \in \cO$, ist auch $\bigcap_{i = 1}^n U_i \in \cO$.
        \item Für jede Familie $\set{U_i}_{i \in I} \in \cO$ ist auch $\bigcup_{i \in I} U_i \in \cO$.
    \end{enumerate}
    Ein \emph{topologischer Raum}\index{topologischer Raum} ist ein Paar $(X, \cO)$ einer Menge $X$ und einer Topologie $\cO$ auf $X$.
\end{defi}

\begin{defi}{Stetigkeit in topologischen Räumen}{continuous}
    Seien $(X, \cO)$ und $(Y, \cO')$ topologische Räume. Dann heißt eine Funktion $f: X \to Y$ \emph{stetig}\index{stetig}, falls für jedes $V \subseteq Y$ offen auch $f^{-1}(V)$ offen in $X$ ist (also falls $V \in \cO' \Rightarrow f^{-1}(V) \in \cO$ gilt).
\end{defi}
Man sagt immer noch \emph{Urbilder offener Mengen sind offen}.

Sei $(X,d)$ ein metrischer Raum. Dann ist $(X, \cO = \set{U \subseteq X \mid \text{$U$ offen bezüglich $d$}})$ ein topologischer Raum. \emph{Begründung:} Folgt aus Lemma~\ref{lem:metricSpaceTraits}.

        $\cO$ ist die \emph{von der Metrik induzierte Topologie}.\index{Topologie!von Metrik induziert}

        Ein topologischer Raum heißt \emph{metrisierbar}\index{topologischer Raum!metrisierbar}, falls die Topologie von einer Metrik induziert wird.

\begin{bsp}\leavevmode\label{bsp:discrete}
    \begin{enumerate}[$(i)$]
        \item Sei $X$ eine Menge. In der \emph{diskreten Topologie}\index{Topologie!diskrete Topologie} auf $X$ sind alle Teilmengen offen. Dies ist induziert durch die diskrete Metrik 
        \begin{equation*}
            d(x,y) := \begin{cases}
                0 & x = y, \\
                1 & \text{sonst.}
            \end{cases}
        \end{equation*}
        \begin{proof}
            Für $x \in X$ ist $\set{x} = U(x, \frac{1}{2})$ offen. Jedes $A \subseteq X$ lässt sich schreiben als $A = \bigcup_{a \in A} \set{a}$, ist also als Vereinigung offener Mengen offen.
        \end{proof}
        Für jedes $t > 0$ definiere die Metrik $d_t$ auf $X$ mit 
        \begin{equation*}
            d(x,y) := \begin{cases}
                0 & x = y, \\
                t & \text{sonst.}
            \end{cases}
        \end{equation*}
        Dann induziert $d_t$ die diskrete Topologie auf $X$.
        \item Die \emph{indiskrete Topologie}\index{Topologie!indiskrete Topologie} auf $X$ ist definiert als $\cO := \set{\emptyset, X}$. Enthält $X$ mindestens zwei Punkte, so ist die indiskrete Topologie nicht metrisierbar.
        \begin{proof}
            Seien $x \neq y \in X$ und $d$ eine Metrik auf $X$. Sei $0 < \varepsilon < d(x,y)$. Dann ist $x \in U(x,\varepsilon)$, aber $y \notin U(x, \varepsilon)$. Also ist $U(x, \varepsilon)$ eine offene Teilmenge ungleich $\emptyset$ von $X$. Also induziert $d$ nicht die indiskrete Topologie. 
        \end{proof}
    \end{enumerate}
\end{bsp}

\begin{bemerkung}
    Sei $X$ eine Menge mit der indiskreten Topologie und $(Y, \cO)$ ein beliebiger topologischer Raum. Dann ist jede Funktion $f: Y \to X$ stetig.
\end{bemerkung}
\begin{proof}
    Sei $f: Y \to X$ eine Funktion. In $X$ sind nur $\emptyset$ und $X$ offen. Es ist $f^{-1}(\emptyset) = \emptyset$ und $f^{-1}(X) = Y$, diese Mengen sind nach Definition offen in $Y$ nach Definition~\ref{def:topology}.
\end{proof}

\begin{defi}{Äquivalenz von Metriken}{metricEquiv}
    Zwei Metriken $d_1, d_2$ auf einer Menge $X$ heißen \emph{äquivalent}\index{Metrik!äquivalent}, falls es $c_1, c_2 > 0$ gibt, sodass
    \begin{equation}
        c_1 \cdot d_1(x,y) \leq d_2(x,y) \leq c_2 \cdot d_1(x,y) \quad \quad \forall x,y \in X. \label{eq:metricEquiv}
    \end{equation}
    Dies ist eine symmetrische Relation,~\eqref{eq:metricEquiv} impliziert 
    \begin{equation*}
        \frac{1}{c_2} d_2(x,y) \leq d_1(x,y) \leq \frac{1}{c_1} d_2(x,y).
    \end{equation*}
\end{defi}

\begin{beispiel}
    Für je zwei $t, t' > 0$ sidn $d_t$ und $d_{t'}$ aus Beispiel~\ref{bsp:discrete} äquivalennt für $c_1 := \frac{t}{t'}, c_2 := \frac{t'}{t}$.
\end{beispiel}

\begin{thm}{Äquivalente Metriken induzieren dieselbe Topologie}{equivalentMetrics}
    Äquivalente Metriken induzieren die gleiche Topologie.
\end{thm}
\begin{proof}
    Wegen der Symmetrie genügt es zu zeigen, dass jede bezüglich $d_2$ offene Menge auch offen bezüglich $d_1$ ist.

    Sei $U \subseteq X$ offen bezüglich $d_2$ und sei $x \in U$. Also existiert ein $\varepsilon' > 0$ mit $U_{d_2}(x, \varepsilon') \subseteq U$.

    Dann ist $U_{d_1}(x, \frac{\varepsilon'}{c_2}) \subseteq U_{d_2}(x, \varepsilon') \subseteq U$, da $d_1(x,y) < \frac{\varepsilon'}{c_2}$ impliziert, dass $d_2(x,y) < c_2 \cdot \frac{\varepsilon'}{c_2} = \varepsilon'$. Also ist $U$ offen bezüglich $d_1$.
\end{proof}

\begin{bemerkung}
    Auch nicht äquiavlente Metriken könen dieselbe Topologie induzieren (vgl. Aufgabe 1.2).
\end{bemerkung}

\begin{thm}{Normen äquivalent, wenn induzierte Topologien gleich}{normInduceTopology}
    Sei $V$ ein reeller Vektorraum, und $\norm{\cdot}_1$ und $\norm{\cdot}_2$ Normen auf $V$. Dann sind äquivalent:
    \begin{enumerate}[$(i)$]
        \item Die Normen $\norm{\cdot}_1$ und $\norm{\cdot}_2$ sind äquivalent,
        \item $\norm{\cdot}_1$ und $\norm{\cdot}_2$ induzieren dieselbe Topologie auf $V$.
    \end{enumerate}
    Auf $\IR^n$ sind je zwei Normen äquivalent, also induzieren je zwei Normen dieselbe Topologie.
\end{thm}

\begin{defi}{Umgebung}{topSurrounding}
    Sei $X$ ein topologischer Raum, $U \subseteq X$ und $x \in U$. Dan heißt $U$ \emph{Umgebung von $x$}\index{Umgebung}, falls eine offene Menge $O \subseteq X$ existiert, sodass $x \in O \subseteq U$.
\end{defi}

\begin{bemerkung}
    Für metrische Räume ist dies äquivalent zur Definition~\ref{def:surrounding}.
\end{bemerkung}

\begin{thm}{Offen $\Leftrightarrow$ Umgebung all seiner Punkte}{openSurrounding}
    Sei $X$ ein topologischer Raum, $U \subseteq X$. Dann ist $U$ genau dann offen in $X$, wenn es Umgebung jedes seiner Punkte ist.
\end{thm}
\begin{proof}
    \begin{itemize}
        \item[\enquote{$\Rightarrow$}] Setze $O = U$ für jedes $x \in U$.
        \item[\enquote{$\Leftarrow$}] Für $x \in U$ wähle $U_x \subseteq X$ offen mit $x \in U_x \subseteq U$. Dann ist 
        $U = \bigcup_{x \in U} U_x$ offen als Vereinigung offener Mengen.
    \end{itemize}
\end{proof}

\begin{defi}{Abgeschlossene Mengen}{closedSet}
    Sei $X$ ein topologischer Raum. Eine Teilmenge $A \subseteq X$ heißt \emph{abgeschlossen}\index{abgeschlossene Menge}, falls ihr Komplement $X \setminus A := \set{x \in X \mid x \notin A}$ offen in $X$ ist.
\end{defi}
\begin{bemerkung}
    Für Teilmengen metrischer Räume stimmt dies mit dem Begriff aus der Analysis überein.\footnote{Oder es ist einfach dieselbe Definition.}
\end{bemerkung}


%%%%%%%%%%%%%%%%%%%%%%%%%%%%%%%%%%%%%%%%%%%%%%%%%%%%%%%%%%%%%%%%%%%%%%%%%%%%%%%%%%%%%%%%%%%%%%%%%%%%
%%%%%%%%%%%%%%%%%%%%%%%%%%%%%%%%%%%%%%%%%% 2024 / 04 / 16 %%%%%%%%%%%%%%%%%%%%%%%%%%%%%%%%%%%%%%%%%%
%%%%%%%%%%%%%%%%%%%%%%%%%%%%%%%%%%%%%%%%%%%%%%%%%%%%%%%%%%%%%%%%%%%%%%%%%%%%%%%%%%%%%%%%%%%%%%%%%%%%

\breakline{Beginn Vorlesung 03 (16.04.2024)}

\begin{bemerkung}
    Ein topologischer Raum lässt sich auch übber die Menge der abgeschlossenen Teilmengen charakterisieren. Diese müssen folgende Axiome erfüllen:
    \begin{itemize}
        \item[$(A1)$] $X$ und $\emptyset$ sind abgeschlossen
        \item[$(A2)$] Für $A_1, \ldots, A_n$ abgeschlossen ist auch $\bigcup_{i = 1}^n A_i$ abgeschlossen
        \item[$(A3)$] Für $\set{A_i}_{i \in I}$ abgeschlossen ist auch $\bigcap_{i \in I} A_i$ abgeschlossen 
    \end{itemize}
\end{bemerkung}
\begin{proof}
    Wir nutzen
    \begin{equation*}
        X \setminus \left(\bigcup_{i = 1}^n A_i\right) = \bigcap_{i = 1}^n \left(X \setminus A_i\right)
    \end{equation*}
    und
    \begin{equation*}
        X \setminus \left(\bigcap_{i \in I} A_i\right) = \bigcap_{i \in I} \left(X \setminus A_i\right)
    \end{equation*}
    und dann folgt die Aussage aus der Definition.
\end{proof}

\begin{lem}{Stetigkeit $\Leftrightarrow$ Urbilder abgeschlossener Mengen sind abgeschlossen}{continuousClosed}
    Seien $X, Y$ topologische Räume und $f: X \to Y$ eine Funktion. Dann ist $f$ genau dann stetig, wenn für alle abgeschlossenen Teilmengen $A \subseteq Y$ die Menge $f^{-1}(A)$ abgeschlossen in $X$ ist.
\end{lem}
\begin{proof}
    Nutze $f^{-1}(Y \setminus A) = X \setminus f^{-1}(A)$ und Definiton~\ref{def:continuous}.
\end{proof}

\begin{beispiel}
    Es folgen einige topologische Räume:
    \begin{enumerate}
        \item \emph{Sierpinski-Raum:}\index{Topologie!Sierpinski-Raum} Sei $X = \set{a,b}$ und $\cO = \set{\emptyset, X, \set{a}}$. Dort hat $b$ nur eine Umgebung, nämlich $X$. Dieser Raum ist nicht metrisierbar.
        \item \emph{Ko-endliche Topologie:}\index{Topologie!ko-endliche} Sei $X = \IN$ und $\cO = \set{U \subseteq \IN \mid \IN \setminus U \text{ ist endlich}} \cap \set{\emptyset}$. Auch dieser Raum ist nicht metrisierbar (Vgl. Aufgaben 1 und 3).
        \item \emph{Zariski-Topologie:}\index{Topologie!Zariski} Sei $X = \IC^n$. Wir nennen eine Menge $A \subseteq \IC^n$ abgeschlossen, falls die Familie $\set{f_i \in \IC[x_1, \ldots, x_n]}_{i \in I}$, also durch den Polynomring, gegeben sind, sodass $A = \set{x \in \IC^n \mid f_i(x) = 0 \text{ für alle $i \in I$}}$.
        
        Wenn $n = 1$, ist dies die ko-endliche Topologie mit $\IC$ statt $\IN$.
    \end{enumerate}
\end{beispiel}


\chapter{Unterräume und Quotientenräume}
Wir stellen uns in diesem Kapitel folgende Fragen:

\textbf{Frage 1.} Sei $Y \subseteq X$.
\begin{enumerate}[$(a)$]
    \item Wenn $X$ eine Metrik hat, erbt $Y$ eine Metrik? -- \emph{Ja, eingeschränkte Metrik.}
    \item Wenn $X$ eine Topologie hat, erbt $Y$ eine Topologie? -- \emph{Ja, Unterraumtopologie.}
\end{enumerate}

\textbf{Frage 2.} Sei $X$ eine Menge mit einer Äquivalenzrelation $\sim$, und sei $Z$ die Menge der Äquivalenzklassen.
\begin{enumerate}[$(a)$]
    \item Wenn $X$ eine Metrik hat, erbt $Z$ eine Metrik? -- \emph{Nein.}
    \item Wenn $X$ eine Topologie hat, erbt $Z$ eine Topologie? -- \emph{Ja, Quotiententopologie.}
\end{enumerate}

\begin{thm}{Unterraumtopologie}{subspaceTopology}
    Sei $(X, \cO)$ ein topologischer Raum und $Y \subseteq X$ eine Teilmenge. Dann definiert 
    \begin{equation*}
        \cO_y := \set{U \subseteq Y \mid \text{Es gibt $U' \subseteq X$ offen, sodass $U = U' \cap Y$}}
    \end{equation*}
    eine Topologie auf $Y$, genannt \emph{Unterraumtopologie}\index{Unterraumtopologie}.
\end{thm}

\begin{proof}
    Wir überprüfen die Axiome einer Topologie.
    \begin{enumerate}[$(i)$]
        \item $\emptyset = \emptyset \cap Y$, $Y = X \cap Y$ sind offen, da $\emptyset$ und $X$ offen in $X$ sind.
        \item Seien $U_1, \ldots, U_n \subseteq Y$ offen. Dan existiert $U_i' \subseteq X$ offen mit $U_i = U_i' \cap Y$. Damit ist 
        \begin{equation*}
            \bigcap_{i = 1}^n U_i = \underset{\text{offen in $X$}}{\underbrace{\bigcap_{i = 1}^n U_i}} \cap Y
        \end{equation*}
        offen in $Y$.
        \item Sei $\set{U_i}_{i \in I}$ eine Familie offener Teilmengen von $Y$. Dann existieren $\set{U_i'}_{i \in I}$ offen in $X$, sodass $U_i = U_i' \cap Y$. Also ist 
        \begin{equation*}
            \bigcup_{i \in I} U_i = \underset{\text{offen in $X$}}{\underbrace{\bigcap_{i \in I} U_i}} \cap Y
        \end{equation*}
        offen in $Y$.
    \end{enumerate}
\end{proof}

\begin{lem}{Universelle Eigenschaft der Unterraumtopologie}{universalTraitSubspace}\index{Universelle Eigenschaft!Unterraumtopologie}
    Sei $X$ ein topologischer Raum und $Y \subseteq X$ offen. Dann ist die Inklusion $i: Y \to X$ stetig.
\end{lem}

\begin{proof}
    Sei $U \subseteq X$ offen. Dann ist $i^{-1}(U) = Y \cap U$ offen nach Definition.
\end{proof}
Das ist die \emph{größte und maximale Topologie auf $Y$}, sodass $i: Y \to X$ stetig ist.

\textbf{Warnung.} Für $Z \subseteq Y \subseteq X$ muss man unterscheiden zwischen \emph{offen in $Y$} und \emph{offen in $X$}. Wenn $Z$ offen in $X$ ist, ist $Z = Z \cap Y$ offen in $Y$. Die Umkehrung gilt im Allgemeinen nicht.

Sei $X = \IR^2$ und $Y = \IR \times \set{0}$ bezüglich der metrischen Topologie. Sei $Z = \intervaloo{0, \infty} \times \set{0}$. Dann ist $Z$ offen in $Y$, weil $Z = \underset{\text{offen in $X$}}{\underbrace{\intervaloo{0,\infty} \times \IR}} \cap Y$, aber $Z$ ist nicht offen in $\IR^2 = X$.

\begin{bemerkung}
    Sei $(X,d)$ ein metrischer Raum, $Y \subseteq X$. Dann ist die Unterraumtopologie auf $Y$ bezüglich der induzierten Topologie auf $X$ gleich der induzierten Topologie bezüglich der eingeschränkten Metrik auf $Y$.
\end{bemerkung}

\begin{proof}
    Sei $y \in Y$. Dann ist $U_{d|_Y}(y, \varepsilon) = U_d(y, \varepsilon) \cap Y$.
\end{proof}

Es gibt zwei Wege, Quotiententopologien zu definieren; über eine Äquivalenzrelation oder über surjektive Abbildungen.

Sei $X$ eine Menge mit einer Äquivalenzrelation $\sim$ (also reflexiv, symmetrisch und transitiv). Für $x \in X$ definieren wir $[x] := \set{x' \in X \mid x \sim x}$ die \emph{Äquivalenzklasse} von $x$. Außerdem definieren wir $\nicefrac{X}{\sim} := \set{[x] \mid x \in X}$ die Menge der Äquivalenzklassen von $\sim$.

Wir definieren 
\begin{align*}
    q: X \surj& \nicefrac{X}{\sim}, \\
    x \mapsto& [x].   
\end{align*}
Die Abbildung $q$ ist surjektiv und heißt \emph{kanonische Projektion}.

Umgekehrt, sei $q: X \surj Z$ surjektiv. Wir definieren $x \sim y :\Leftrightarrow q(x) = q(y)$. Dies ist eine Äquivalenzrelation auf $X$, und 
\begin{align*}
    \nicefrac{X}{\sim} \bij& Z, \\
    [x] \mapsto & q(x)
\end{align*}
ist bijektiv.

Damit bestimmen sich Äquivalenzrelationen und surjektive Abbildungen gegenseitig.

\begin{thm}{Quotiententopoloie}{quotientTopology}
    Sei $X$ ein topologischer Raum mit Äquivalenzrelation $\sim$ und kanonischer Projektion $q: X \to \nicefrac{X}{\sim} = Z$. Dann definiert 
    \begin{equation*}
        \cO_Z := \set{U \subseteq Z \mid q^{-1}(U) \text{ ist offen in }X}
    \end{equation*}
    eine Topologie auf $Z$, die \emph{Quotiententopologie}\index{Quotiententopologie}.
\end{thm}

\begin{proof}
    Wir überprüfen die Axiome einer Topologie.
    \begin{enumerate}[$(i)$]
        \item $q^{-1}(\emptyset) = \emptyset$ und $q^{-1}(Z) = X$ sind offen in $X$, damit sind $\emptyset$ und $Z$ offen in $Z$.
        \item Seien $U_1, \ldots, U_n \subseteq Z$ offen. Dann ist 
        \begin{equation*}
            q^{-1}\left(\bigcap_{i = 1}^n U_i\right) = \bigcap_{i = 1}^n q^{-1}(U_i)
        \end{equation*}
        offen in $X$, da jedes $q^{-1}(U_i)$ nach Voraussetzung offen in $X$ ist. Also ist $\bigcap_{i = 1}^n U_i$ offen in $Z$.
        \item Sei $\set{U_i}_{i \in I}$ eine Familie offener Teilmengen von $Z$. Dann ist 
        \begin{equation*}
            q^{-1}\left(\bigcup_{i \in I} U_i\right) = \bigcup_{i \in I} \underset{\text{offen in $X$}}{\underbrace{q^{-1}(U_i)}}
        \end{equation*}
        offen in $X$. Also ist $\bigcup_{i \in I} U_i$ offen in $Z$.
    \end{enumerate}
\end{proof}

\begin{bemerkung}
    Die Projektion $q: X \to \nicefrac{X}{\sim} = Z$ ist stetig und die Quotiententopologie ist maximal unter allen Topologien auf $Z$, für die $q$ stetig ist.
\end{bemerkung}

\begin{thm}{Universelle Eigenschaft der Quotiententopologie}{universalTraitQuotientTopo}\index{Universelle Eigenschaft!Quotiententopologie}
    Sei $f: X \to Y$ eine stetige Abbildung zwischen topologischen Räumen und sei $q: X \to Z$ die Projektion zu einer Äquivalenzrelation auf $X$. Angenommen, es gibt eine (notwendigerweise eindeutige) Abbildung $g: Z \to Y$ mit $f = g \circ q$. Dann ist $g$ stetig, das folgende Diagramm kommutiert also. 
    \[
        \begin{tikzcd}
            X \arrow[rr, "f"] \arrow[dr, twoheadrightarrow, "q"]   &                           & Y \\
                                                & Z \arrow[ur, dotted, "g"]
        \end{tikzcd}
    \]
\end{thm}

\begin{proof}
    Sei $U \subseteq Y$ offen. Wir wollen zeigen, dass $g^{-1}(U)$ offen in $Z$ ist. Es ist
    \begin{equation*}
        q^{-1}(g^{-1}(U)) = (g \circ q)^{-1}(U) = f^{-1}(U)
    \end{equation*}
    offen in $X$, da $f$ stetig ist. 

    Nach Definition der Quotiententopologie ist also $g^{-1}(U)$ offen in $Z$, also ist $g$ stetig.
\end{proof}

\begin{beispiel}
    Sei $X = \intervalcc{0,1}$ mit der metrischen Topologie und Äquivalenzrelation erzeugt von $0 \sim 1$. \emph{Ziel:} $\nicefrac{X}{\sim}$ \enquote{ist ein Kreis}.
\end{beispiel}

\begin{defi}{Homöomorphie}{homeomorphy}
    Zwei topologische Räume $X, Y$ heißen \emph{topologisch äquivalent} bzw. \emph{homöomorph}\index{topologische Äquivalenz}\index{Homöomorphie}, falls stetige Funktionen $f: X \to Y$ und $g: Y \to X$ existieren, sodass $g \circ f = \id_X$ und $f \circ g = \id_Y$.
\end{defi}

\begin{thm}{$\nicefrac{\intervalcc{0,1}}{\sim} \iso \IS^1$}{circleEquivLine}
    Der Quotientenraum $\nicefrac{\intervalcc{0,1}}{\sim}$ ist topologisch äquivalent zum Kreis $\IS^1 = \set{x \in \IR^2 \mid \norm{x} = 1}$.
\end{thm}

\begin{proof}
    Betrachte $f' : \intervalcc{0,1} \to \IS^1 \subseteq \IC$ mit $f'(t) = e^{2\pi i t} = (\cos(2\pi t), \sin(2\pi t))$. Diese Funktion ist stetig nach Ana I. 
    
    Es ist $f'(0) = 1 = f'(1)$, also existiert die Funktion $f: \nicefrac{\intervalcc{0,1}}{\sim} \to \IS^1$ mit $[x] \mapsto f'(x)$. Nach Satz~\ref{thm:universalTraitQuotientTopo} ist $f$ stetig, das Diagramm 
    \[
        \begin{tikzcd}
            \intervalcc{0,1} \arrow[rr, "f"] \arrow[dr, twoheadrightarrow, "q"]   &                           & \IS^1 \\
                                                & \nicefrac{\intervalcc{0,1}}{\sim} \arrow[ur, dotted, "g"]
        \end{tikzcd}
    \]
    kommutiert also. Weiter ist $f$ bijektiv, also existiert eine Umkehrfunktion $g: \IS^1 \to \nicefrac{\intervalcc{0,1}}{\sim}$. Es bleibt zu zeigen, dass $g$ stetig ist. \emph{Äquivalent:} Für alle $U \subseteq \nicefrac{\intervalcc{0,1}}{\sim}$ offen ist $f(U)$ offen in $\IS^1$ (da $g^{-1}(U) = f(U))$. Man sagt auch, $f$ ist eine \emph{offene Abbildung}\index{offene Abbildung}.

    Sei $U \subseteq \nicefrac{\intervalcc{0,1}}{\sim}$ offen, $x \in U$. Wir zeigen, dass $f(U)$ Umgebung von $f(x)$ ist.

    \emph{Fall 1: $x \neq \set{0,1} = q(0) = q(1)$.} Dann ist $x = \set{t} = q(t)$ für ein $0 < t < 1$. Das Urbild $q^{-1}(U) \subseteq \intervalcc{0,1}$ ist offen und enthält $t$. Somit existiert ein $\varepsilon > 0$ mit $\intervaloo{t-\varepsilon, t+\varepsilon} \subseteq q^{-1}(U)$.
    
    Also enthält $f'(q^{-1}(U)) = f(U)$ ein offenes Kreisbogenstück um $f'(t) = f(x)$. Somit ist $f(U)$ Umgebung von $f(x)$.

    \emph{Fall 2: $x = \set{0,1}$.} Das Urbild $q^{-1}(U) \subseteq \intervalcc{0,1}$ ist offen und enthält $0$ und $1$. Sei $\varepsilon > 0$ mit $\intervalco{0,\varepsilon} \cup \intervaloc{1-\varepsilon, 1} \subseteq q^{-1}(U)$. Dann enthält $f(U) = f'(q^{-1}(U))$ ein offenes Kreisbogenstück um $1 = f(x)$. Somit ist $f(V)$ Umgebung von $f(x)$.

    Damit ist die Aussage gezeigt und $g$ ist stetig.
\end{proof}


%%%%%%%%%%%%%%%%%%%%%%%%%%%%%%%%%%%%%%%%%%%%%%%%%%%%%%%%%%%%%%%%%%%%%%%%%%%%%%%%%%%%%%%%%%%%%%%%%%%%
%%%%%%%%%%%%%%%%%%%%%%%%%%%%%%%%%%%%%%%%%% 2024 / 04 / 18 %%%%%%%%%%%%%%%%%%%%%%%%%%%%%%%%%%%%%%%%%%
%%%%%%%%%%%%%%%%%%%%%%%%%%%%%%%%%%%%%%%%%%%%%%%%%%%%%%%%%%%%%%%%%%%%%%%%%%%%%%%%%%%%%%%%%%%%%%%%%%%%

\breakline{Beginn Vorlesung 04 (18.04.2024)}

\textbf{Warnung.} Stetige, bijektive Abbildungen müssen im Allgemeinen keine topologischen Äquivalenzen sein.

\begin{beispiel}\leavevmode
    \begin{enumerate}[$(i)$]
        \item Sei $X$ eine Menge. Dann ist die Identität
        \begin{equation*}
            \begin{tikzcd}
                (X, \text{diskrete Topologie}) \arrow[r, "\id"] & (X, \text{indiskrete Topologie})
            \end{tikzcd}
        \end{equation*}
        stetig und bijektiv, aber keine topologische Äquivalenz, falls $X$ mindestens zwei Elemente besitzt.
        \item Wir betrachten
        \begin{align*}
            \exp: \intervaloo{0,1} \bij& \IS^1 \\
            t \mapsto& e^{2\pi i t}.
        \end{align*}
        Diese Abbildung ist stetig und bijektiv, aber nicht offen, denn $\intervalco{0, 1/2} \subseteq \intervalco{0,1}$ ist offen, aber $\exp\left(\intervalco{0,1/2}\right) \subseteq \IS^1$ ist nicht offen.
    \end{enumerate}
\end{beispiel}

\begin{beispiel}
    Wir betrachten einige weiteren Quotientenräume.
    \begin{enumerate}[$(1)$]
        \item Sei $X = \intervalcc{0,1}^2 \subseteq \IR^2$ mit $\sim$ erzeugt von 
        \begin{itemize}
            \item $(t,0) \sim (t,1)$ für alle $t \in \intervalcc{0,1}$ und
            \item $(0,s) \sim (1, 1-s)$ für alle $s \in \intervalcc{0,1}$.
        \end{itemize}
        Dabei wird ein Quadrat in die \emph{Kleinsche Flasche}\index{Kleinsche Flasche} transformiert.
        \item Sei $X = \IR^{n+1} \setminus \set{0}$ mit $\sim$ erzeugt von $x \sim \lambda x$ für alle $x \in X$ und $\lambda \in \IR \setminus \set{0}$. Damit entspricht ein Punkt in $\nicefrac{X}{\sim}$ einer Gerade in $\IR^{n+1} \setminus \set{0}$. Diesen Raum nennt man \emph{Projektiver Raum}\index{Projektiver Raum} $\IR P^n$.
        \item Sei $X$ ein topologischer Raum und $A \subseteq X$. Wir definieren $\sim$ auf $X$ als Äquivalenzrelation erzeugt durh $a \sim a'$ für alle $a, a' \in A$. Wir schreiben $\nicefrac{X}{A}$ für $\nicefrac{X}{\sim}$. 
        
        Es ist $\nicefrac{\intervalcc{0,1}}{\set{0,1}} \iso \IS^1$. Für $D^n := \set{x \in \IR^n \mid \abs{x} \leq 1}$, also die \emph{$n$-dimensionale Scheibe}, gilt $\nicefrac{D^n}{\IS^{n-1}} \iso \IS^n$ (vgl. Aufgabe 8).
        \item Der Raum $\nicefrac{\intervalcc{0,1}}{\intervaloc{0,1}}$ hat zwei  Punkte: $\set{0}$ (abgeschlossen) und $\intervaloc{0,1}$ (nicht abgeschlossen). Dies ist der \emph{Sierpinski-Raum}\index{Sierpinski-Raum} (nicht metrisierbar).
    \end{enumerate}
\end{beispiel}

\chapter{Trennungsaxiome und Kompaktheit}
\begin{defi}{Hausdorff-Raum}{hausdorff}
    Ein topologischer Raum $X$ heißt \emph{Hausdorff-Raum} (oder er ist \emph{hausdorffsch})\index{Hausdorff-Raum}, wenn sich je zwei verschiedene Punkte durch offene Mengen trennen lassen: Für $x \neq y \in X$ existieren offene Mengen $U_x$ und $U_y$, sodass $x \in U_x, y \in  U_y$ und $U_x \cap U_y = \emptyset$. Diese Eigenschaft heißt auch \emph{Trennungsaxiom $T2$}\index{Trennungsaxiom}\index{Trennungsaxiom!T2}.
\end{defi}

\begin{thm}{Metrische Topologie ist Hausdorff}{metricHausdorff}
    Die induzierte Topologie einer Metrik ist immer hausdorffsch.
\end{thm}
\begin{proof}
    Sei $(X,d)$ ein metrischer Raum, $x \neq y \in X$. Setze $U_x := U(x, d(x,y)/2)$ und $U_y := U(y, d(x,y)/2)$. Dann ist $x \in U_x, y \in U_y$. Weiter ist $U_x \cap U_y = \emptyset$: Angenommen, $z \in U_x \cap U_y$. Dann ist $d(x,y) \leq d(x,z) + d(z,y) < d(x,y)/2 + d(x,y)/2 = d(x,y) ~\lightning$.
\end{proof}

\begin{bemerkung}
    \begin{enumerate}[$(i)$]
        \item Sei $X$ ein Hausdorff-Raum, $A \subseteq X$. Dann ist $A$ auch ein Hausdorff-Raum.
        \begin{proof}
            Sei $x \neq y \in A$. Seien $U_x', U_y' \subseteq X$ offen mit $x \in U_x', y \in U_y'$ und $U_x' \cap U_y' = \emptyset$. Dann ist $x \in U_x := U_x' \cap A$ und $y \in U_y := U_y' \cap A$. Es folgt $U_x \cap U_y \subseteq U_x' \cap U_y' = \emptyset$.
        \end{proof}
        \item Sei $X$ eine Menge mit $\geq 2$ Elementen. Dann ist die induzierte Topologie auf $X$ nicht hausdorffsch.
        \item Jede diskrete Topologie ist hausdorffsch.
        \begin{proof}
            Setze $U_x := \set{x}, U_y := \set{y}$ für $x \neq y$. Dann ist $\set{x} \cap \set{y} = \emptyset$.
        \end{proof}
        \item Jeder endliche Hausdorff-Raum trägt die diskrete Topologie.
        \begin{proof}
            Sei $X$ endlich und hausdorffsch, $x \in X$. Für jedes $x \neq y \in X$ wähle $U_x^{(y)}$ offen mit $x \in U_x^{(y)}$ und $y \notin U_x^{(y)}$. Dann ist 
            \begin{equation*}
                \set{x} = \bigcap_{y \in X \setminus \set{x}} U_x^{(y)}
            \end{equation*}
            offen als endlicher Schnitt über offene Teilmengen. Also ist die Topologie auf $X$ diskret.
        \end{proof}
    \end{enumerate}
\end{bemerkung}

\textbf{Erinnerung.} Für $X \subseteq \IR^n$ sind äquivalent:
\begin{enumerate}[$(i)$]
    \item $X$ ist abgeschlossen und beschränkt,
    \item Jede Folge in $X$ hat eine konvergente Teilfolge,
    \item Jede offene Überdeckung hat eine endliche Teilüberdeckung.
\end{enumerate}
$X$ heißt \emph{kompakt}\index{Kompaktheit}, wenn $(i)$-$(iii)$ erfüllt sind.

\begin{defi}{Quasikompaktheit, Kompaktheit}{quasiCompact}
    Ein topologischer Raum $X$ heißt \emph{quasikompakt}\index{Quasikompaktheit}, falls jede offene Überdeckung eine endliche Teilüberdeckung besitzt:

    Für jede Familie offener Teilmengen $\set{U_i}_{i \in I}$ von $X$ mit 
    \begin{equation*}
        X = \bigcup_{i \in I} U_i
    \end{equation*}
    existieren $i_1, \ldots, i_n$, sodass
    \begin{equation*}
        X = \bigcup_{k = 1}^n U_{i_k}.
    \end{equation*}
    $X$ heißt \emph{kompakt}\index{Kompaktheit}, falls es quasikompakt und hausdorffsch ist.
\end{defi}

\begin{beispiel}
    Jeder endliche, topologische Raum ist quasikompakt.
\end{beispiel}

\begin{thm}{$\intervalcc{0,1}$ kompakt}{compact01}
    Das Intervall $\intervalcc{0,1}$ ist kompakt.
\end{thm}
\begin{proof}
    Die \emph{Hausdorffeigenschaft} folgt aus der Metrik.

    Zur \emph{Quasikompaktheit}: Sei $\set{U_i}_{i \in I}$ eine offene Überdeckung von $\intervalcc{0,1}$. Betrachte 
    \begin{equation*}
        X := \set{t \in \intervalcc{0,1} \mid \text{es gibt endliche Teilüberdeckung von $\intervalcc{0,t}$}}.
    \end{equation*}
    Setze $s := \sup(X)$. (Dieses existiert, da $X$ nichtleer $(0 \in X)$ und beschränkt ist). Wir zeigen, dass $s \in X$:
    
    Wähle dazu $i \in I$ mit $s \in U_i$. Sei $\varepsilon > 0$, sodass $\intervaloo{s - \varepsilon, s + \varepsilon} \subseteq U_i$. Da $s$ das Supremum ist, gibt es ein $t \in \intervaloc{s - \varepsilon, s}$, sodass $t \in X$. Damit gibt es $i_1, \ldots, i_n \in I$, sodass
    \begin{equation*}
        \intervalcc{0,t} \subseteq \bigcup_{k = 1}^n U_{i_k}
    \end{equation*}
    und damit ist $s \in X$.

    Nun zeigen wir, dass $s = 1$. Sei wieder $i \in I$ mit $s \in U_i$ und sei $\varepsilon > 0$, sodass $\intervaloo{s - \varepsilon, s + \varepsilon} \subseteq U$. 
    
    Angenommen, $s < 1$. Da $s \in X$, gibt es $i_1, \ldots, i_n$ mit $\intervalcc{0,s} \subseteq \bigcup_{i = 1}^n U_i$. Sei $t' \in \intervaloo{s, s+\varepsilon}$. Dann ist
    \begin{equation*}
        \intervalcc{0,t'} \subseteq \bigcup_{k = 1}^n U_{i_k}, \quad \text{aber $t' \in X$}. 
    \end{equation*}
    Das ist ein Widerspruch zu $s = \sup(X) ~\lightning.$
\end{proof}

\begin{thm}{Abgeschlossene Teilmenge quasikompakter Räume ist quasikompakt}{quasiCompactSubspace}
    Sei $X$ quasikompakt und $Y \subseteq X$ abgeschlossen. Dann ist $Y$ quasikompakt.
\end{thm}
\begin{proof}
    Sei $\set{U_i}_{i \in I}$ eine offene Überdeckung von $Y$. Seien $U_i' \subseteq X$ offen mit $U_i' \cap Y = U_i$ für alle $i \in I$. Dann ist $\set{U_i'} \cup \set{X \setminus Y}$ eine offene Überdeckung von $X$. 

    Da $X$ quasikompakt ist, gibt es $i_1, \ldots, i_n \in I$, sodass 
    \begin{equation*}
        X = \bigcup_{k = 1}^n U_{i_k}' \cup (X \setminus Y).
    \end{equation*}
    Dann ist
    \begin{equation*}
        Y = Y \cap X = \bigcup_{k = 1}^n \underset{=U_{i_k}}{\underbrace{(Y \cap U_{i_k}')}} \cup \underset{= \emptyset}{\underbrace{(Y \cap (X \setminus Y))}}.
    \end{equation*}
    Damit ist $\bigcup_{k = 1}^n U_{i_k}$ eine endliche Teilüberdeckung von $Y$.
\end{proof}

\begin{thm}{Quasikompakter Unterraum von Hausdorff ist abgeschlossen}{HausdorffClosedSubspace}
    Sei $X$ ein Hausdorff-Raum und $Y \subseteq X$ ein quasikompakter Unterraum von $X$. Dann ist $Y$ abgeschlossen in $X$.
\end{thm}

\textbf{Korollar.} Sei $X$ kompakt und $Y \subseteq X$. Dann ist $Y$ genau dann kompakt, wenn $Y$ abgeschlossen in $X$ ist.
\begin{proof}
    Folgt direkt aus Satz~\ref{thm:quasiCompactSubspace} und Satz~\ref{thm:HausdorffClosedSubspace}.
\end{proof}

\begin{lem}{Hausdorfftrennung zwischen Punkt und Menge}{HausdorffPointSet}
    Seien $X, Y$ wie in Satz~\ref{thm:HausdorffClosedSubspace}. Dann gibt es für jeden Punkt $x \in X \setminus Y$ offene Teilmengen $U_{x,Y}$ und $U_{x,Y}'$, sodass $x \in U_{x,Y}$, $Y \subseteq U_{x,Y}'$ und $U_{x,Y} \cap U_{x,Y}' = \emptyset$.
\end{lem}

\begin{proof}
    Für jedes $y \in Y$ existieren offene Teilmengen $U_{x,y}, U_{x,y}'$, sodass $x \in U_{x,y}, y \in U_{x,y}'$ mit $U_{x,y} \cap U_{x,y}' = \emptyset$, da $X$ ein Hausdorffraum ist. Dann ist $\set{U_{x,y}' \cap Y}_{y \in Y}$ eine offene Überdeckung von $Y$.
    
    Da $Y$ quasikompakt ist, gibt es $y_1, \ldots, y_n \in Y$, sodass $Y \subseteq \bigcup_{k = 1}^n U_{x,y_k}'$.

    Wir setzen $U_{x,Y}' := \bigcup_{k = 1}^n U_{x,y_k}'$ und $U_{x,Y} := \bigcap_{k = 1}^n U_{x,y}$.

    Dann ist $x \in U_{x,Y}$ und $Y \subseteq U_{x,Y}'$ und $U_{x,Y} \cap U_{x,y_i}' \subseteq U_{x,y_i} \cap U_{x,y_i}' = \emptyset$. Damit folgt auch $U_{x,Y} \cap U_{x,Y}' = \emptyset$.
\end{proof}

\begin{proof}[Beweis: Satz~$\ref{thm:HausdorffClosedSubspace}$]
    Für jedes $x \in X \setminus Y$ gibt es nach Lemma~\ref{lem:HausdorffPointSet} eine offene Teilmenge $U_{x,Y}$ mit $x \in U_{x,Y}$und $Y \cap U_{x,Y} = \emptyset$. Also ist
    \begin{equation*}
        X \setminus Y = \bigcup_{x \in X \setminus Y} U_{x,Y}
    \end{equation*} 
    offen und somit $Y$ abgeschlossen.
\end{proof}

\begin{beispiel}\leavevmode
    \begin{enumerate}[$(i)$]
        \item $\intervalco{0,1} \subseteq \intervalcc{0,1}$ ist nicht quasikompakt. 
        \item Wenn $X$ indiskret ist, und $Y \subseteq X$ nicht abgeschlossen ist, sind sowohl $X$ als auch $Y$ quasikompakt, aber $Y$ ist nicht abgeschlossen.
    \end{enumerate}
    
\end{beispiel}



%%%%%%%%%%%%%%%%%%%%%%%%%%%%%%%%%%%%%%%%%%%%%%%%%%%%%%%%%%%%%%%%%%%%%%%%%%%%%%%%%%%%%%%%%%%%%%%%%%%%
%%%%%%%%%%%%%%%%%%%%%%%%%%%%%%%%%%%%%%%%%% 2024 / 04 / 23 %%%%%%%%%%%%%%%%%%%%%%%%%%%%%%%%%%%%%%%%%%
%%%%%%%%%%%%%%%%%%%%%%%%%%%%%%%%%%%%%%%%%%%%%%%%%%%%%%%%%%%%%%%%%%%%%%%%%%%%%%%%%%%%%%%%%%%%%%%%%%%%

\breakline{Beginn Vorlesung 05 (23.04.2024)}

\begin{thm}{Quasikompaktheit unter stetiger Surjektion bleibt erhalten}{surjectiveQuasicompact}
    Sei $f: X \to Y$ stetig und surjektiv. Falls $X$ quasikompakt ist, dann ist auch $Y$ quasikompakt.
\end{thm}
\begin{proof}
    Sei $\set{U_i}_{i \in I}$ eine offene Überdeckung von $Y$. Dann ist $\set{f^{-1}(U_i)}_{i \in I}$ eine offene Überdeckung von $X$, da $f$ stetig ist.

    $X$ ist quasikompakt, also gibt es $i_1, \ldots, i_n \in I$, sodass
    \begin{equation*}
        X = \bigcup_{k = 1}^n f^{-1}(U_{i_k}).
    \end{equation*}
    Da $f$ surjektiv ist, gilt $f(f^{-1}(U_i)) = U_i$ für alle $i \in I$. Also folgt 
    \begin{equation*}
        Y = f(x) = \bigcup_{k = 1}^n f(f^{-1}(U_{i_k})) = \bigcup_{k = 1}^n U_i.
    \end{equation*}
    Also ist $U_{i_1}, \ldots, U_{i_n}$ eine endliche Teilüberdeckung von $Y$.
\end{proof}

\textbf{Korollar 1.} Sei $f: X \to Y$ stetig, $X$ quasikompakt und $Y$ hausdorffsch. Dann ist $f$ eine \emph{abgeschlossene Abbildung}\index{abgeschlossene Abbildung}, d.h. für $A \subseteq X$ abgeschlossen ist auch $f(A) \subseteq Y$ abgeschlossen. 

\begin{proof}
    Es ist $A \subseteq X$ abgeschlossen. Dann ist nach Satz~\ref{thm:quasiCompactSubspace} auhch $A$ quasikompakt. Nach Satz~\ref{thm:surjectiveQuasicompact} auch $f(A)$ quasikompakt und mit Satz~\ref{thm:HausdorffClosedSubspace} auch $f(A) \subseteq Y$ abgeschlossen.
\end{proof}

\textbf{Korollar 2.} Sei $f: X \to Y$ stetig, bijektiv, $X$ quasikompakt, $Y$ hausdorffsch. Dann ist $f$ eine topologische Äquivalenz.

\begin{proof}
    Wir müssen zeigen, dass $f^{-1}$ stetig ist, also dass für alle $A \subseteq X$ abgeschlossen ist $\left(f^{-1}\right)^{-1}(A) = f(A) \subseteq Y$ abgeschlossen ist. Das folgt direkt aus Korollar 1.
\end{proof}

\textbf{Korollar 3.} Sei $f: X \to Y$ stetig und surjektiv. Seien $X$ quasikompakt und $Y$ hausdorffsch. Dann trägt $Y$ automatisch die Quotiententopologie bezüglich $f$, also $U \subseteq Y$ ist genau dann offen, wenn $f^{-1}(U) \subseteq X$ offen ist.

\begin{proof}\leavevmode
    \begin{itemize}
        \item[\enquote{$\Rightarrow$}] folgt direkt aus der Stetigkeit von $f$.
        \item[\enquote{$\Leftarrow$}] Sei $f^{-1}(U) \subseteq X$ offen. Dann ist $X \setminus f^{-1}(U)$ abgeschlossen, also nach Korollar 1 auch $f(X \setminus f^{-1}(U)) = Y \setminus U$ abgeschlossen (nach Surjektivität von $f$). Also ist $U$ offen in $Y$.
    \end{itemize}
\end{proof}

\begin{beispiel}[zu Korollar 2]
    Die Funktion $\exp: \intervalcc{0,1} \to \IS^1, \exp(t) := e^{2\pi i t}$ ist stetig und surjektiv, faktorisiert über stetige, bijektive Abbildung $f: \nicefrac{\intervalcc{0,1}}{0 \sim 1} \to \IS^1$. 

    Da $\intervalcc{0,1}$ quasikompat ist, ist auch $\nicefrac{\intervalcc{0,1}}{0 \sim 1}$ quasikompakt. $\IS^1$ trägt die metrische Topologie, ist also ein Hausdorff-Raum. Mit Korollar 2 folgt dann schon die topologische Äquivalenz.
\end{beispiel}

\begin{defi}{Normaler topologischer Raum}{normalTopology}
    Ein topologischer Raum $X$ heißt \emph{normal}\index{normal}\index{topologischer Raum!normal}, falls
    \begin{enumerate}[$(i)$]
        \item er hausdorffsch ist,
        \item für je $A, B \subseteq X$ abgeschlossen mit $A \cap B = \emptyset$ offene Teilmengen $U_A, U_B \subseteq X$ existieren mit $A \subseteq U_A, B \subseteq U_B$ und $U_A \cap U_B = \emptyset$.
    \end{enumerate}
    Es wird auch als \emph{Trennungsaxiom $T4$}\index{Trennungsaxiom}\index{Trennungsaxiom!T4} bezeichnet.
\end{defi}

\begin{thm}{Kompakt $\Rightarrow$ normal}{compactSpaceNormal}
    Jeder kompakte Raum ist normal.
\end{thm}
\begin{proof}
    Seien $A, B \subseteq X$ abgeschlossen und disjunkt. Da $X$ kompakt ist, und $B$ abgeschlossen ist, ist $B$ quasikompakt (nach Satz~\ref{thm:quasiCompactSubspace}). Dann folgt aus Lemma~\ref{lem:HausdorffPointSet}, dass für jedes $a \in A$ Mengen $U_a, U_a' \subseteq X$ offen existieren mit $a \in U_a, B \subseteq U_a'$ und $U_a \cap U_a' = \emptyset$. 
    Dann ist $\set{A \cap U_a}_{a \in A}$ eine offene Überdeckung von $A$. Da $A$ als abgeschlossene Menge in $X$ kompakt ist, gibt es $a_1, \ldots, a_n \in A$, sodass
    \begin{equation*}
        A \subseteq \bigcup_{k = 1}^n U_{a_k}.
    \end{equation*}
    Wir setzen
    \begin{equation*}
        U_A := \bigcup_{k = 1}^n U_{a_k} \quad \quad \text{und} \quad \quad U_B := \bigcap_{k = 1}^n U_{a_k}'.
    \end{equation*}
    Dann ist $A \subseteq U_A, B \subseteq U_B$ mit $U_A, U_B$ offen und $U_{a_i} \cap U_B \subseteq U_{a_i} \cap U_{a_i}' = \emptyset$. Also ist dann auch $U_A \cap U_B = \emptyset$.
\end{proof}

\begin{thm}{Metrischer Raum $\Rightarrow$ normal}{metricSpaceNormal}
    Jeder metrische Raum ist normal.
\end{thm}
\begin{proof}
    Sei $X$ ein metrischer Raum. Nach Satz~\ref{thm:metricHausdorff} ist $X$ hausdorffsch.

    Seien $A, B \subseteq X$ abgeschlossen und disjunkt. Da $B$ abgeschlossen ist, gibt es für alle $a \in A$ ein $\varepsilon_a > 0$, sodass $U(a, \varepsilon_a) \cap B = \emptyset$. Analog gibt es für jedes $b \in B$ ein $\varepsilon_b > 0$, sodass $U(b, \varepsilon_b) \cap A = \emptyset$.

    Setze 
    \begin{equation*}
        U_A := \bigcup_{a \in A} U\left(a, \frac{\varepsilon_a}{2}\right) \quad \quad \text{und} \quad \quad U_B := \bigcup_{b \in B} U\left(b, \frac{\varepsilon_b}{2}\right),
    \end{equation*} 
    beide Mengen sind als Vereinigung offener Mengen wieder offen. Dann ist $A \subseteq U_A, B \subseteq U_B$. Weiter ist dabei für $a \in A$ und $b \in B$ 
    \begin{equation*}
        U\left(a, \frac{\varepsilon_a}{2}\right) \cap U\left(b, \frac{\varepsilon_b}{2}\right) = \emptyset:
    \end{equation*}
    Gäbe es ein $z \in U\left(a, \frac{\varepsilon_a}{2}\right) \cap U\left(b, \frac{\varepsilon_b}{2}\right)$, so folgt 
    \begin{equation*}
        d(a,b) \leq d(a,z) + d(b,z) < \frac{\varepsilon_a}{2} + \frac{\varepsilon_b}{2}.
    \end{equation*}
    Andererseits ist $d(a,b) \geq \varepsilon_a$ und $d(a,b) \geq \varepsilon_b$, also ist 
    \begin{equation*}
        d(a,b) = \frac{d(a,b)}{2} + \frac{d(a,b)}{2} \geq \frac{\varepsilon_a}{2} + \frac{\varepsilon_b}{2}
    \end{equation*}
    und somit kann $z$ nicht existieren. $\lightning$
\end{proof}


\chapter{Das Urysohnsche Lemma und der Fortsetzungssatz von Tietze}

\begin{thm}{Urysohnsches Lemma}{urysohn}
    Sei $X$ ein normaler Raum, $A, B \subseteq X$ abgeschlossen und disjunkt. Dann existiert eine stetige Funktion $f: X \to \intervalcc{0,1}$, sodass $f(a) = 0$ für alle $a \in A$ sowie $f(b) = 1$ für alle $b \in B$.
\end{thm}

\begin{bemerkung}
    Im Allgemeinen ist es nicht möglich, $f$ so zu wählen, dass $A = f^{-1}(0)$ und $B = f^{-1}(1)$.
\end{bemerkung}

\textbf{Vorbereitung.} Sei $X$ ein topologischer Raum und $Y \subseteq X$ eine Teilmenge. Dann heißt 
\begin{equation*}
    \overline{Y} := \bigcap_{\substack{A \supseteq Y \\ \text{abg.}}} A 
\end{equation*}
der \emph{Abschluss}\index{Abschluss (Menge)} von $Y$ sowie
\begin{equation*}
    \mathring{Y} := \bigcup_{\substack{B \subseteq Y \\ \text{offen}}} U
\end{equation*}
das \emph{Innere}\index{Inneres (Menge)} von $Y$. Dabei ist $\overline{Y}$ abgeschlossen und $\mathring{Y}$ offen.

\begin{lem}{Zwischenmenge im normalen Raum}{normalSetBetween}
    Sei $X$ ein normaler Raum, $A \subseteq X$ abgeschlossen und $U \subseteq X$ offen mit $A \subseteq U$. Dann existiert $V \subseteq X$ offen mit 
    \begin{equation*}
        A \subseteq V \subseteq \overline{V} \subseteq U.
    \end{equation*}
\end{lem}

\begin{proof}
    Setze $B := X \setminus U$ abgeschlossen und disjunkt von $A$. 

    Da $X$ normal ist, gibt es $U_1, U_2$ offen mit $A \subseteq U_1$ und $B \subseteq U_2$ sowie $U_1 \cap U_2 = \emptyset$. Wir setzen $V := U_1$. Dann ist 
    \begin{equation*}
        \overline{V} = \overline{U_1} \subseteq X \setminus U_2 \subseteq X \setminus B = U,
    \end{equation*}
    wobei die erste Inklusion folgt, weil $X \setminus U_2$ abgeschlossen ist und $U_1 \subseteq X \setminus U_2$.
\end{proof}

\begin{proof}[Beweis von Satz~$\ref{thm:urysohn}$]
    Für jedes $p \in \intervalcc{0,1} \cap \IQ$ konstruieren wir eine offene Teilmenge $U_p \subseteq X$, sodass $\overline{U_p} \subseteq U_q$ für $p < q$. Hierfür wählen wir eine Abzählung $p_1, p_2, \ldots$ von $\intervalcc{0,1} \cap \IQ$, sodass $p_1 = 1, p_2 = 0$ und induktiv definieren wir $U_{p_i}$ folgendermaßen:
    
    Für $U_{p_1} = U_1 = X \setminus B$.

    Für $U_{p_2} = U_0$ verwenden wir Lemma~\ref{lem:normalSetBetween} mit $U = X \setminus B$ und erhalten $V \subseteq X$ offen mit $A \subseteq V \subseteq \overline{V} \subseteq X \setminus B$. Wir setzen $U_0 := V$.

    Für $n \geq 3$ nehmen wir an, dass $U_{p_1}, \ldots, U_{p_{n-1}}$ bereits definiert sind mit $\overline{U_{p_i}} \subseteq U_{p_j}$ für $p_i < p_j$.

    Da die Menge $\set{p_1, \ldots, p_n}$ wohlgeordnet ist, gibt es einen direktn Vorgänger $p_i$ von $p_n$ und einen direkten Nachfolger $p_j$ von $p_n$. Dann ist $p_i < p_j$, also ist $\overline{U_{p_i}} \subseteq U_{p_j}$. Nach Lemma~\ref{lem:normalSetBetween} existiert $V \subseteq X$ offen mit
    \begin{equation*}
        \overline{U_{p_i}} \subseteq V \subseteq \overline{V} \subseteq U_{p_j}.
    \end{equation*}
    Wir setzen $U_{p_n} := V$ und erweitern die Definition von $U_p$ auf alle $p \in \IQ$ mit $U_p = \emptyset$ für $p < 0$ und $U_p = X$ für $p > 1$.

    Für $x \in X$ setzen wir
    \begin{equation*}
        \IQ(x) := \set{p \in \IQ \mid x \in U_p} \subseteq \IQ_{\geq 0}.
    \end{equation*}
    diese Menge ist nichtleer und nach unten beschränkt. Damit setzen wir $f(x) := \inf (\IQ(x)) \in \intervalcc{0,1}$. Diese Funktion erfüllt die gewünschten Eigenschaften:

    Für $a \in A$ ist $\IQ(a) = \intervalco{0,a} \cap \IQ$, also ist $f(a) = 0$.

    Für $b \in B$ ist $\IQ(b) = \intervaloo{1,\infty} \cap \IQ$, da $b \notin U_1 = X \setminus B$, also ist $f(b) = 1$.

    Für die Stetigikeit von $f$ stellen wir fest: Wenn $x \in \overline{U_p}$, dann ist $f(x) \leq p$ (da $x \in U_q$ für $q > p$). Analog ist $f(x) \geq p$ für $x \notin \overline{U_p}$.

    Sei $U \subseteq \IR$ offen und $x \in f^{-1}(U)$. Dann gibt es $a < b \in \IR$, sodass $\intervaloo{a,b} \subseteq U$ und $a < f(x) < b$. Wir wählen $p, q \in \IQ$ mit $a < p < f(x) < q < b$. 
    
    \emph{Behauptung:} Dann enthält die offene Menge $U_q \setminus \overline{U_p}$ schon $x$ und ist enthalten in $f^{-1}(U)$.

    Sei $x \in U_q \setminus \overline{U_p}$. Für $y \in U_q \setminus \overline{U_p}$ ist $f(y) \leq q$ und $f(y) \geq p$, aber $f(y) \in \intervaloo{a,b} \subseteq U$. Damit ist $y \in f^{-1}(U)$ und $f$ ist stetig.
\end{proof}

%%%%%%%%%%%%%%%%%%%%%%%%%%%%%%%%%%%%%%%%%%%%%%%%%%%%%%%%%%%%%%%%%%%%%%%%%%%%%%%%%%%%%%%%%%%%%%%%%%%%
%%%%%%%%%%%%%%%%%%%%%%%%%%%%%%%%%%%%%%%%%% 2024 / 04 / 25 %%%%%%%%%%%%%%%%%%%%%%%%%%%%%%%%%%%%%%%%%%
%%%%%%%%%%%%%%%%%%%%%%%%%%%%%%%%%%%%%%%%%%%%%%%%%%%%%%%%%%%%%%%%%%%%%%%%%%%%%%%%%%%%%%%%%%%%%%%%%%%%

\breakline{Beginn Vorlesung 06 (25.04.2024), heute mit Qi Zhu\footnote{\enquote{Wie ihr vielleicht merkt, ist Markus heute ein bisschen geschrumpft.} -- Qi Zhu, 25.04.2024}}

\begin{thm}{Fortsetzungssatz von Tietze}{tietzeContinuation}
    Sei $X$ ein normaler Raum und $A \subseteq X$ abgeschlossen.
    \begin{enumerate}[$(i)$]
        \item Jede stetige Funktion $f: A \to \intervalcc{0,1}$ lässt sich zu einer stetigen Funktion $f: X \to \intervalcc{0,1}$ fortsetzen.
        \item Jede stetige Funktion $f: A \to \IR$ lässt sich zu einer stetigen Funktion $f: X \to \IR$ fortsetzen.
    \end{enumerate}
\end{thm}

\begin{bem}
    Urysohns Lemma (Satz~\ref{thm:urysohn}) ist ein Spezialfall von Tietze.

    \begin{proof}
        Seien $A, B \subseteq X$. Betrachte $A \sqcup B \to \intervalcc{0,1}$ mit $a \mapsto 0$ und $b \mapsto 1$ für alle $a \in A, b \in B$. Damit lässt sich diese Funktion mit Satz~\ref{thm:tietzeContinuation} zu $X \to \intervalcc{0,1}$ fortsetzen.
    \end{proof}
\end{bem}

\begin{proof}[Beweis von Tietze $(\ref{thm:tietzeContinuation})$]\leavevmode
    \begin{enumerate}[$(i)$]
        \item Wir konstruieren eine Folge stetiger Funktionen $s_n: X \to \intervalcc{0,1}$, sodass
        \begin{enumerate}[$(a)$]
            \item $(s_n)_{n \in \IN}$ gleichmäßig gegen eine Funktion $s$ konvergiert,
            \item $s$ sich auf $A$ zu $f$ einschränkt.
        \end{enumerate}
        Wir erinnern kurz daran, was \emph{gleichmäßige Konvergenz} ist:
        \begin{enumerate}[$(1)$]
            \item $s_n \to s$ \emph{gleichmäßig} $\Leftrightarrow \forall \varepsilon>0 \exists N \in IN: \forall x \in X, n \geq \IN: d(s_n(x), s(x)) < \varepsilon$,
            \item $s_n$ sind alle stetig $\Rightarrow s$ ist stetig.
        \end{enumerate}
        Wir beweisen zunächst $(2)$: Sei $U \subseteq \intervalcc{0,1}$ offen und sei $x \in s^{-1}(U)$. Damit gibt es ein $\varepsilon > 0$, sdass $U(s(x), \varepsilon) \subseteq U$.

        Da $s_n \to s$ gleichmäßig, gibt es ein $N \in \IN$, sodass für alle $x' \in X, n \geq N$ gilt:
        \begin{equation*}
            d(s_n(x'), s(x')) < \frac{\varepsilon}{3}.
        \end{equation*}
        Da $s_N$ stetig ist, gibt es $V \subseteq X$ offen mit $x \in V$, sodass 
        \begin{equation*}
            s(V) \subseteq U\left(s_N(x), \frac{\varepsilon}{3}\right).
        \end{equation*}
        Damit gilt für alle $\tilde{x} \in V$:
        \begin{equation*}
            d(s(x), s(\tilde{x})) \leq d(s(x), s_N(x)) + d(s_N(x), s_N(\tilde{x})) + d(s_N(\tilde{x}), s(\tilde{x})) < \frac{\varepsilon}{3} + \frac{\varepsilon}{3} + \frac{\varepsilon}{3} = \varepsilon.
        \end{equation*}
        Damit ist $s(V) \subseteq U(s(x), \varepsilon) \subseteq U$ und $V \subseteq s^{-1}(U)$ ist eine offene Umgebung von $x$, also ist $s$ stetig.

        Wir basteln uns eine Hilfskonstruktion: \\
        \emph{Input:} $\alpha: A \to \intervalcc{-r, r}$ stetig für ein $r \in \intervalco{0,\infty}$,\\
        \emph{Output:} $g: X \to \intervalcc{-\frac{1}{3} r, \frac{1}{3} r}$.

        Wir setzen $I_1 := \intervalcc{-r, -\frac{1}{3}r}, I_2 := \intervalcc{-\frac{1}{3}r, \frac{1}{3}r, I_3 := \intervalcc{\frac{1}{3}r, r}}$. Sei $B := \alpha^{-1}(I_1) \subseteq A$ abgeschlossen und $C := \alpha^{-1}(I_3) \subseteq A$ abgeschlossen. Dann gibt es nach Satz~\ref{thm:urysohn} eine stetige Funktion $g: X \to \intervalcc{-\frac{1}{3}r, \frac{1}{3}r}$ mit 
        \begin{equation*}
            g(b) = -\frac{1}{3} r \quad \forall b \in B \quad \quad \text{und} \quad \quad g(c) = \frac{1}{3}r \quad \forall c \in C.
        \end{equation*}
        Dieses $g$ erfüllt $\abs{g(a) - \alpha(a)} \leq \frac{2}{3}r$ für alle $a \in A$, da:
        
        \emph{Fall 1:} $a \in B \Rightarrow \alpha(a) \in I_1, g(a) = -\frac{1}{3}r$,

        \emph{Fall 2:} $a \in C \rightarrow \alpha(a) \in I_3, g(a) = \frac{1}{3}r$,

        \emph{Fall 3:} $a \in A \setminus (B \cup C) \Rightarrow \alpha(a) \in I_2, g(a) \in I_2$.

        Wir ersetzen oBdA $\intervalcc{0,1}$ durch $\intervalcc{-1,1}$, also $f: A \to \intervalcc{-1,1}$ (wende ggf. eine lineare Transformation an).

        \emph{Konstruktion:}
        \begin{enumerate}[1.]
            \item Hilfskonstruktion mit $\alpha = f, r = 1 \Rightarrow$ wir erhalten $g_1: X \to \intervalcc{\frac{-1}{3}, \frac{1}{3}}$ stetig mit $\abs{g_1(a) - f(a)} \leq \frac{2}{3}$. Wir setzen nun $s_1 = g_1: X \to \IR$.
            \item Hilfskonstrukion mit $\alpha := f - s_1|_A: A \to \intervalcc{-\frac{2}{3}, \frac{2}{3}}$, also $r = \frac{2}{3}$. Damit erahlten wir $g_2: X \to \intervalcc{-\frac{1}{3} \cdot \frac{2}{3}, \frac{1}{3} \cdot \frac{2}{3}}$ stetig mit 
            \begin{equation*}
                \abs{g_1(a) + g_2(a) - f(a)} = \abs{g_2(a) - (f(a) - s_1(a))} \leq \frac{2}{3} \cdot \frac{2}{3}.
            \end{equation*}
            Wir setzen nun $s_1 := g_1 + g_2: X \to \IR$.
            \item Hilfskonstrukion mit $\alpha := f - s_2|_A : A \to \intervalcc{-\left(\frac{2}{3}\right)^2, \left(\frac{2}{3}\right)^2}$ und $r = \left(\frac{2}{3}\right)^2$ usw.
        \end{enumerate}
        Dies setzen wir induktiv fort und erhalten stetige Funktionen $g_1, \ldots, g_n$ auf $X$ mit 
        \begin{align*}
            g_n: X \to \intervalcc{-\frac{1}{3} \cdot \left(\frac{2}{3}\right)^{n-1}, \frac{1}{3} \cdot \left(\frac{2}{3}\right)^{n-1}}, \\
            \abs{(g_1(a) + \cdots + g_n(a)) - f(a)} \leq \left(\frac{2}{3}\right)^n,
        \end{align*}
        und wir setzen $s_n := g_1 + \cdots + g_n: X \to \IR$.

        \emph{Hilfsrechnung:} Für $k > n$ ist 
        \begin{align*}
            \abs{s_k(x) - s_n(x)} &= \abs{\sum_{i = n + 1}^k g_i(x)} \overset{\triangle\text{-UG}}{\leq} \sum_{i = n+1}^k \abs{g_i(x)} \\
            &\leq \frac{1}{3}\sum_{i = n + 1}^k \left(\frac{2}{3}\right)^i < \frac{1}{3} \sum_{i = n+1}^\infty \left(\frac{2}{3}\right)^i\\  
            &= \frac{1}{3} \left(\frac{2}{3}\right)^n \sum_{j = 0}^\infty \left(\frac{2}{3}\right)^j \\
            \overset{\text{geoReihe}}&{=} \frac{1}{3} \left(\frac{2}{3}\right)^n \frac{1}{1- \frac{2}{3}} = \left(\frac{2}{3}\right)^n
        \end{align*}
        \emph{Folgerung 1.} Für jedes $x \in X$ ist $(s_(x))_n$ eine Cauchy-Folge in $\IR$, also gibt es $s(x) \in \IR$, sodas $s_n(x) \to s(x)$.

        \emph{Folgerung 2.} Es gilt 
        \begin{equation*}
            \abs{s_n(x) - s(x)} = \lim_{k \to \infty} \abs{s_k(x) - s_n(x)} \leq \left(\frac{2}{3}\right)^n,
        \end{equation*}
        also konvergiert $s_n \to s$ gleichmäßig, und $s: X \to \IR$ ist stetig.

        Es bleibt zu zeigen, dass 
        \begin{enumerate}[1.]
            \item $s ~ f$ fortsetzt,
            \item das Bild von $s$ in $\intervalcc{-1,1}$ liegt.
        \end{enumerate}
        \emph{Zu $1.$:} Es gilt 
        \begin{equation*}
            \abs{s_n(a) - f(a)} = \abs{(g_1(a) + \cdots + g_n(a)) - f(a)} \leq \left(\frac{2}{3}\right)^n
        \end{equation*}
        für $a \in A$. Für $n \to \infty$ folgt $\abs{s(a) - f(a)} = 0$, also $s(a) = f(a)$.
        
        \emph{Zu $2.$:} Wir können es entweder direkt überprüfen, oder wie folgt:

        Ersetze $s$ durch $s' = r \circ s$ mit 
        \begin{equation*}
            r: \IR \to \intervalcc{-1,1}, \quad x \mapsto \begin{cases}
                x & x \in \intervalcc{-1,1}, \\
                \frac{x}{\abs{x}} & \text{sonst.}
            \end{cases}
        \end{equation*}
        Damit gilt die Aussage und $(i)$ ist bewiesen.
        
        \item Sei $f: A \to \IR$ stetig. OBdA ersetzen wir $\IR$ durch das topologisch äquivalente $\intervaloo{0,1}$, verfahren also ähnlich für $f: A \to \intervaloo{-1,1}$ stattdessen.
        
        \emph{Hilfskonstruktion:} Nach $(i)$ existiert ein stetiges $s: X \to \intervalcc{-1,1}$ mit $s|_A = f$. Wir setzen $D := s^{-1}(\set{-1}) \cup s^{-1}(\set{1})$, also ist $D$ abgeschlossen in $X$.

        Da $A \cap D = \emptyset$ ($s(A) = f(A) \subseteq \intervaloo{-1,1}$), gibt es nach Satz~\ref{thm:urysohn} eine stetige Funktion $\varphi: X \to \intervalcc{0,1}$ mit $\varphi(D) = \set{0}$ und $\varphi(A) = \set{1}$.

        \emph{Konstruktion:} Setze 
        \begin{equation*}
            h: X \to \intervaloo{-1,1}, \quad x \mapsto \varphi(x) \cdot s(x).
        \end{equation*}
        Es bleibt zu zeigen:
        \begin{enumerate}[$(a)$]
            \item Das Bild von $h$ liegt in $\intervaloo{-1,1}$,
            \item $h$ setzt $f$ fort.
        \end{enumerate}

        \emph{Zu $(a)$:} Es gibt zwei Fälle.

        \emph{Fall 1:} $x \in D \Rightarrow h(x) = \varphi(x) \cdot s(x) = 0 \in \intervaloo{-1,1}$, da $\varphi(D) = \set{0}$.

        \emph{Fall 2:} $x \notin D \Rightarrow \abs{h(x)} = \underset{<1}{\underbrace{\abs{\varphi(x)}}} \cdot \underset{\leq 1}{\underbrace{\abs{s(x)}}} < 1$.

        \emph{Zu $(b)$:} Es ist $h(a) = \varphi(a) \cdot s(a) = s(a) = f(a)$.
    \end{enumerate}
    Damit ist der Beweis abgeschlossen. Puh!
\end{proof}

\section{Anwendung: Pseudokompaktheit in metrischen Räumen}

\begin{prop}
    Ein metrischer Raum $(X,d)$ ist genau dann (quasi-)kompakt, wenn jede stetige Funktion $f: X \to \IR$ beschränkt ist. \emph{Man sagt, dass $X$ pseudokompakt ist.}\index{pseudokompakt}
\end{prop}

\begin{proof}
    \begin{itemize}
        \item[\enquote{$\Rightarrow$}] Da $X$ kompakt ist, ist $f(X)$ kompakt, also beschränkt.
        \item[\enquote{$\Leftarrow$}] Angenommen, $X$ ist nicht kompakt. Dann gibt es eine Folge $(x_n)_{n \in \IN}$ in $X$ ohne konvergente Teilfolge. Damit sind alle $x_n$ paarweise verschieden und jede konvergente Folge in $A = \set{x_n \mid n \in \IN}$ ist irgendwann konstant. 
        Daraus folgt, dass alle Folgen in $A$, die in $X$ konvergieren, einen Grenzwert in $A$ haben. Damit ist $A \subseteq X$ abgeschlossen und diskret (dasselbe Argument für jede Teilmenge aus $A$ zeigt dies mit unendlichem Abstieg).
        
        Definiere $f: A \to \IR, x_n \mapsto n$. Diese Funktion ist also stetig (Satz~\ref{lem:continuousClosed}). Daraus folgt mit Satz~\ref{thm:tietzeContinuation}, dass wir $f$ erweitern können zu $\tilde{f}: X \to \IR$ mit $\tilde{f}|_A = f$. Diese Funktion ist unbeschränkt.\footnote{\enquote{Eigentlich wollte ich noch ein Beispiel zur Cantor-Menge machen, aber Markus wollte, dass ich lieber weitermache.} -- Qi Zhu, 25.04.2024}
    \end{itemize}
\end{proof}


\section{Quotientenräume und Hausdorff-Eigenschaft}
\begin{defi}{Saturation}{saturation}
    Seien $X, Z$ Mengen und $q: X \to Z$ eine Funktion. Eine Teilmenge $W \subseteq X$ \emph{saturiert}\index{Saturation}, falls 
    \begin{equation*}
        W = q^{-1}(q(W)) \Leftrightarrow \forall x \in X \setminus W: ~ q(x) \notin q(W).
    \end{equation*}
\end{defi}

\begin{lem}{Eigenschaften der Saturation}{saturationProperties}
    Sei $q: X \to Z$ eine Funktion und $A \subseteq X$ saturiert. Dann ist 
    \begin{enumerate}[$(i)$]
        \item $q(X \setminus A) = q(X) \setminus q(A)$, 
        \item $X \setminus A$ saturiert.
    \end{enumerate}
\end{lem}

\begin{proof}
    \begin{enumerate}[$(i)$]
        \item \begin{itemize}
            \item[\enquote{$\supseteq$}:] gilt immer,
            \item[\enquote{$\subseteq$}:] Sei $z \in q(X \setminus A)$ geschrieben als $z = q(x)$ mit $x \in X \setminus A$.
            
            Angenommen, $q(x) \in q(A)$. Dann folgt $x \in q^{-1}(q(A)) = A$, da $A$ saturiert ist. $\lightning$
        \end{itemize}
        \item Es gilt 
        \begin{equation*}
            q^{-1}(q(X \setminus A)) \overset{(i)}{=} q^{-1}(q(X) \setminus q(A)) = q^{-1}(q(X)) \setminus q^{-1}(q(A)).
        \end{equation*}
    \end{enumerate}
\end{proof}

\begin{lem}{Quotientenraumprojektion offener, saturierter Mengen}{saturationQuotientProjection}
    Sei $q: X \to Z$ eine Quotientenraumprojektion und $U \subseteq X$ offen und saturiert. Dann ist $q(U) \subseteq Z$ offen.
\end{lem}

\begin{proof}
    Da $q^{-1}(q(U)) = U$ offen in $X$ ist, ist $q(U) \subseteq Z$ offen laut der Quotiententopologie (auf $Z$).
\end{proof}

\begin{thm}{Quotientenraum Hausdorff $\Leftrightarrow$ Projektion abgeschlossen}{hausdorffClosedMap}
    Sei $X$ ein kompakter, topologischer Raum und $q: X \to Z$ eine Quotientenraumprojektion. Dann ist $Z$ genau dann hausdorffsch, wenn $q$ eine abgeschlossene Abbildung ist, also $q^{-1}(q(A)) \subseteq X$ für alle $A \subseteq X$ abgeschlossen ist.
\end{thm}

%%%%%%%%%%%%%%%%%%%%%%%%%%%%%%%%%%%%%%%%%%%%%%%%%%%%%%%%%%%%%%%%%%%%%%%%%%%%%%%%%%%%%%%%%%%%%%%%%%%%
%%%%%%%%%%%%%%%%%%%%%%%%%%%%%%%%%%%%%%%%%% 2024 / 04 / 30 %%%%%%%%%%%%%%%%%%%%%%%%%%%%%%%%%%%%%%%%%%
%%%%%%%%%%%%%%%%%%%%%%%%%%%%%%%%%%%%%%%%%%%%%%%%%%%%%%%%%%%%%%%%%%%%%%%%%%%%%%%%%%%%%%%%%%%%%%%%%%%%

\breakline{Beginn Vorlesung 07 (30.04.2024)}

\begin{proof}
    \begin{itemize}
        \item[\enquote{$\Rightarrow$}:] Korollar nach Satz~\ref{thm:HausdorffClosedSubspace}.
        \item[\enquote{$\Leftarrow$}:] \emph{Behauptung 1.} Für ale $z \in Z$ ist $\set{z}$ abgeschlossen.

        \emph{Beweis der Behauptung 1.} Da $q$ surjektiv ist, gibt es $x \in X$ mit $q(x) = z$. Da $X$ ein Hausdorff-Raum ist, ist $\set{x}$ abgeschlossen in $X$. Also ist nach Voraussetzung $q(\set{x}) = \set{z}$ abgeschlossen.

        Seien nun $z_1 \neq z_2 \in Z$. Dann sind nach Behauptung 1 die Mengen $A_1 := q^{-1}(z_1)$ und $A_2 := q^{-1}(z_2)$ abgeschlossen in $X$ und disjunkt.

        Da $X$ kompakt ist, ist $X$ nach Satz~\ref{thm:compactSpaceNormal} auch normal, also gibt es $U_1, U_2 \subseteq X$ offen und disjunkt mit $A_i \subseteq U_i$.

        Setze
        \begin{equation*}
            V_1 := X \setminus q^{-1}\big(q(X \setminus U_1)\big) \quad \text{und} \quad V_2 := X \setminus q^{-1}\big(q(X \setminus U_2)\big).
        \end{equation*}
        Es ist $X \setminus U_i$ abgeschlossen, also ist auch $q^{-1}\bigl(q(X \setminus U_i)\bigr)$ abgeschlossen, also $V_1, V_2$ offen.

        Weiter sind $V_1, V_2$ saturiert und disjunkt mit $A_1 \subseteq V_1$ und $A_2 \subseteq V_2$. Damit sind auch $q(V_1), q(V_2) \subseteq Z$ offen und disjunkt mit $z_i \in q(V_i)$ nach Lemma~\ref{lem:saturationQuotientProjection}.
    \end{itemize}
\end{proof}

\begin{thm}{Q-Raum Hausdorff $\Leftrightarrow$ Q-Menge abgeschlossen, Raum kompakt}{quotientHausdorffClosed}
    Sei $X$ kompakt, $A \subseteq X$. Dann ist $\nicefrac{X}{A}$ genau dann ein Hausdorff-Raum, wenn $A \subseteq X$ abgeschlossen ist.
\end{thm}

\begin{proof}
    \begin{itemize}
        \item[\enquote{$\Rightarrow$}:] Wenn $\nicefrac{X}{A}$ hausdorffsch ist, ist $\set{A} \subseteq \nicefrac{X}{A}$ abgeschlossen. Mit Stetigkeit von $q$ folgt, dass $A = q^{-1}(\set{A})$ abgeschlossen ist.
        \item[\enquote{$\Leftarrow$}:] Nach Satz~\ref{thm:hausdorffClosedMap} genügt es zu zeigen, dass $q: X \to \nicefrac{X}{A}$ abgeschlossen ist.
        
        Sei dazu $B \subseteq X$ abgeschlossen. Dann ist
        \begin{equation*}
            q^{-1}\bigl(q(B)\bigr) = \begin{cases}
                B & A \cap B = \emptyset \\
                A \cup B & \text{sonst}
            \end{cases}
        \end{equation*}
        abgeschlossen. Somit ist auch $q(B) \subseteq \nicefrac{X}{A}$ abgeschlossen.
    \end{itemize}
\end{proof}


\chapter{Basis, Subbasis und Produkttopologie}

\begin{defi}{Basis, Subbasis}{baseSubbase}
    Sei $X$ ein topologischer Raum. Eine Menge offener Mengen $\cS \subseteq \cO$ (für eine Topologie $\cO$) heißt 
    \begin{enumerate}[$(i)$]
        \item \emph{Basis}\index{Basis}, falls jede offene Teilmenge von $X$ Vereinigung von Mengen aus $\cS$ ist.
        \item \emph{Subbasis}\index{Basis!Subbasis}\index{Subbasis}, fals jede offene Teilmenge von $X$ Vereinigung von endlichen Schnitten von Mengen in $\cS$ ist, also ist jedes $U \subseteq X$ offen von der Form 
        \begin{equation*}
            U = \bigcup_{i \in I} \left(\bigcap_{k \in K_i} S_k\right)
        \end{equation*}
        mit $K_i$ endlich und $S_k \in \cS$.
    \end{enumerate}
\end{defi}
\begin{bemerkung}
    Jede Basis ist auch eine Subbasis.
\end{bemerkung}
\begin{beispiel}
    Sei $(X,d)$ ein metrischer Raum. Dann ist $\cS := \set{U(x, \varepsilon) \mid x \in X}$ eine Basis der induzierten Topologie auf $X$.
\end{beispiel}

\begin{lem}{Stetigkeit in der Subbasis}{continuousSubbase}
    Sei $f: X \to Y$ eine Funktion zwischen topologischen Räumen $X$ und $Y$, und $\cS$ eine Subbasis der Topologie auf $Y$. Dann ist $f$ genau dann stetig, wenn für alle $U \in \cS$ die Menge $f^{-1}(U) \subseteq X$ offen ist.
\end{lem}

\begin{proof}
    \begin{itemize}
        \item[\enquote{$\Rightarrow$}:] Folgt aus der Definition von Stetigkeit.
        \item[\enquote{$\Leftarrow$}:] Jedes $U \subseteq Y$ offen ist von der Form 
        \begin{equation*}
            U = \bigcup_{i \in I} \left(\bigcap_{k \in K_i} S_k\right)
        \end{equation*}
        mit $S_k \in \cS$ und $K_i$ endlich. Dann ist 
        \begin{equation*}
            f^{-1}(U) = \bigcup_{i \in I}\left(\bigcap_{k \in K_i} \underset{\text{offen (Vor.)}}{\underbrace{f^{-1}(S_k)}}\right)
        \end{equation*}
        Da $X$ ein topologischer Raum ist, ist dies erneut offen, also ist $f$ stetig.
    \end{itemize}
\end{proof}
\emph{Umgekehrt:} Sei $X$ eine Menge und $\cS$ eine Menge von Teilmengen von $X$. Dann existiert genau eine Topologie $\cO$ auf $X$, für die $\cS$ eine Subbasis ist (die \emph{Topologie erzeugt von $\cS$}):

\begin{equation*}
    \cO := \set{U \subseteq X \mid U = \bigcup_{i \in I} \left(\bigcap_{k \in K_i} \cS_k\right), \text{$K_i$ endlich, $\cS_k \in \cS$}}
\end{equation*}
\begin{proof}
    Übungsaufgabe.
\end{proof}

\begin{thm}{Basis aus Subbasis}{baseOutOfSubbase}
    Seien $X, \cS, \cO$ wie oben. Dann ist $\cS$ genau dann eine Basis von $\cO$, wenn für alle $S_1, S_2 \in \cS$ die Menge $S_1 \cap S_2$ als Vereinigung von Elementen aus $\cS$ geschrieben werden kann.

    \textit{Äquivalent:} Für alle $x \in S_1 \cap S_2$ gibt es $S_3 \in \cS$ mit $x \in S_3 \subseteq S_1 \cap S_2$.
\end{thm}

\begin{proof}
    \begin{itemize}
        \item[\enquote{$\Rightarrow$}:] Da $\cO$ eine Topologie ist, mit $S_1, S_2 \in \cO$, ist $S_1 \cap S_2 \in \cO$. Da $\cS$ eine Basis von $\cO$ ist, ist $S_1 \cap S_2$ Vereinigung von Elementen aus $\cS$.
        \item[\enquote{$\Leftarrow$}:] Jedes $U \subseteq X$ offen ist von der Form 
        \begin{equation*}
            U = \bigcup_{i \in I} \underset{\text{Verein. von El. aus $\cS$}}{\underbrace{\left(\bigcap_{k \in K_i} S_k\right)}}
        \end{equation*}
        mit $S_k \in \cS$ und $K_i$ endlich, also ist $U$ Vereinigung von Elementen aus $\cS$.
    \end{itemize}
\end{proof}

Seien $X_1, X_2$ topologische Räume. \emph{Ziel:} Natürliche Topologie auf 
\begin{equation*}
    X_1 \times X_2 = \set{(x_1, x_2) \mid x_1 \in X_1, x_2 \in X_2}.
\end{equation*}

\emph{Motivation:} $(X_1, d_1), (X_2, d_2)$ metrische Räume. Definiere die \emph{Maximumsmetrik} auf $X_1 \times X_2$ durch 
\begin{equation*}
    d_{\max} \bigl((x_1,x_2),(y_1,y_2)\bigr) = \max\set{d_1(x_1,y_1), d_2(x_2,y_2)}
\end{equation*}
oder alternativ eine äquivalente Metrik 
\begin{equation*}
    d\bigl((x_1,x_2), (y_1,y_2)\bigr) = \sqrt{d_1(x_1,y_1)^2 + d_2(x_2,y_2)^2}.
\end{equation*}

$\varepsilon$-Bälle in der Maximumsmetrik sind \emph{Blöcke}:
\begin{equation*}
    U_{d_{\max}}\bigl((x_1,x_2), \varepsilon\bigr) := U_{d_1}(x_1, \varepsilon) \times U_{d_2}(x_2, \varepsilon).
\end{equation*}

\begin{defi}{Produkttopologie (endlich)}{productTopology}
    Seien $X_1, X_2$ topologische Räume. Die \emph{Produkttopologie}\index{Topologie!Produkt-}\index{Produkttopologie} auf $X_1 \times X_2$ ist die Topologie erzegugt von der Familie 
    \begin{equation*}
        \cB := \set{U_1 \times U_2 \mid U_i \subseteq X_i \text{ offen}}.
    \end{equation*}
\end{defi}

\begin{bemerkung}
    Die Menge $\cB$ ist eine Basis der Produkttopologie, d.h. jede offene Teilmenge $U$ von $X_1 \times X_2$ ist von der Form
    \begin{equation*}
        U = \bigcup_{i \in I} U_1^{(i)} \times U_2^{(i)},
    \end{equation*}
    wobei $U_k^{(i)} \subset X_k$ offen ist.
\end{bemerkung}
\begin{proof}
    \emph{Zu zeigen:} Für $U_1 \times U_2, V_1 \times V_2 \in \cB$ ist $(U_1 \times U_2) \cap (V_1 \times V_2)$ Vereinigung von Elementen aus $\cB$: Es ist 
    \begin{equation*}
        (U_1 \times U_2) \cap (V_1 \times V_2) = \underset{\text{offen in $X_1$}}{\underbrace{U_1 \cap V_1}} \times \underset{\text{offen in $X_2$}}{\underbrace{U_2 \cap V_2}}
    \end{equation*}
    und die Aussage stimmt.
\end{proof}

\begin{beispiel}\leavevmode
    \begin{enumerate}[$(i)$]
        \item Die Standardtopologie auf $\IR^2$ ist die Produkttopologie bezüglich der Standardtopologie auf $\IR$.
        \item $\IR^{n+1} \setminus \set{0}$ ist topologisch äquiavlent zu $\IS^n \times \intervaloo{0,\infty}$.

        Betrachte dazu die Abbildung 
        \begin{align*}
            \varphi: \IR^{n+1} \setminus \set{0} \bij& \IS^n \times \intervaloo{0,\infty} \\
            x \mapsto& \left(\frac{x}{\norm{x}},\norm{x}\right),
        \end{align*}
        $\varphi$ ist bijektiv und stetig mit stetiger Umkehrfunktion
        \begin{align*}
            \psi: \IS^n \times \intervaloo{0,\infty} \bij& \IR^{n+1} \setminus \set{0}, \\
            (y,t) \mapsto& t \cdot y.
        \end{align*}
        \item $\IS^1 \times \IS^1$ ist ein Torus.
        
        Betrachte dazu die Abbildung 
        \begin{align*}
            \varphi: \intervalcc{0,1} \times \intervalcc{0,1} \surj& \IS^1 \times \IS^1 \\
            (s,t) \mapsto& (\exp(2\pi is), \exp(2\pi it)).
        \end{align*}
        $\varphi$ ist stetig und surjektiv und erfüllt $\varphi(s,0) = \varphi(s,1)$ sowie $\varphi(0,t) = \varphi(1,t)$ für alle $s,t \in \intervalcc{0,1}$.

        Also faktorisiert $\varphi$ über 
        \begin{equation*}
            \varphi' : \nicefrac{\intervalcc{0,1} \times \intervalcc{0,1}}{\substack{(s,0) \sim (s,1) \\ (0,t) \sim (1,t) \\ \forall s,t \in \intervalcc{0,1}}} \to \IS^1 \times \IS^1.
        \end{equation*}
        $\varphi'$ ist stetig und bijekti, und $\intervalcc{0,1} \times \intervalcc{0,1}$ ist quasikompakt, damit ist auch $\nicefrac{\intervalcc{0,1} \times \intervalcc{0,1}}{\sim}$ quasikompakt.

        Da $\IS^1 \times \IS^1$ die metrische Topologie trägt, ist $\IS^1 \times \IS^1$ hausdorffsch, also folgt mit Korollar 2 nach Satz~\ref{thm:surjectiveQuasicompact}, dass $\varphi'$ eine topologische Äquivalenz ist.
    \end{enumerate}
\end{beispiel}

\begin{thm}{Universelle Eigenschaft der Produkttopologie}{universalTraitProductTopo}\index{Universelle Eigenschaft!Produkttopologie}
    Seien $A, X_1, X_2$ topologische Räume mit $f_1: A \to X_1, f_2: A \to X_2$ Funktionen. Dann ist die Funktion 
    \begin{align*}
        f: A \to& X_1 \times X_2 \\
        a \mapsto& (f_1(a), f_2(a))
    \end{align*}
    genau dann stetig, wenn $f_1$ und $f_2$ stetig sind.
\end{thm}
\begin{proof}
    \begin{itemize}
        \item[\enquote{$\Rightarrow$}:] Setze
        \begin{align*}
            p_i: X_1 \times X_2 &\to X_i \\
            (x_1, x_2) &\mapsto x_i
        \end{align*}
        (die \emph{kanonischen Projektionen}\index{kanonische Projektion}). Dann ist $f_i = p_i \circ f$ (mit $i \in \set{1,2}$). Da Verknüpfungen stetiger Fuktionen stetig sind, genügt es zu zeigen, dass $p_i$ stetig ist.

        Sei dafür $U_1 \subseteq X_1$ offen. Dann ist 
        \begin{equation*}
            p_1^{-1}(U_1) = U_1 \times X_2.
        \end{equation*}
        Dies ist ein \enquote{Block} in der Definition der Produkttopologie, also offen. Damit ist $p_1$ stetig (für $p_2$ analog), und damit sind auch $f_1$ und $f_2$ stetig.
        \item[\enquote{$\Leftarrow$}:] Seien $f_1$ und $f_2$ stetig. Es genügt zu zeigen, dass $f^{-1}(U_1 \times U_2)$ offen in $A$ ist für alle Basismengen $U_1 \times U_2 \subseteq X_1 \times X_2$ (nach Lemma~\ref{lem:continuousSubbase}).
        
        Es ist 
        \begin{equation*}
            f^{-1}(U_1 \times U_2) = f_1^{-1}(U_1) \cap f_2^{-1}(U_2)
        \end{equation*}
        offen, da $f_i$ stetig sind. Damit ist auch $f$ stetig.
    \end{itemize}
\end{proof}

\begin{thm}{Hausdorff und Quasikompakt unter Produkttopologie}{HausdorffQuasicompactProduct}
    Seien $X, Y$ topologische Räume.
    \begin{enumerate}[$(i)$]
        \item Falls $X$ und $Y$ hausdorffsch sind, dann auch $X \times Y$.
        \item Falls $X$ und $Y$ quasikompakt sind, dann auch $X \times Y$.
    \end{enumerate}
\end{thm}

\begin{proof}
    \begin{enumerate}[$(i)$:]
        \item Seien $(x,y) \neq (x',y') \in X \times Y$. Dann ist $x \neq x'$ oder $y \neq y'$. Falls $x \neq x'$, gibt es $U, U' \subseteq X$ offen mit $x \in U, x' \in U'$ und $U \cap U' = \emptyset$. Dann sind $U times Y$ und $U' \times Y$ offen, $(x,y) \in U \times Y$ und $(x',y) \in U' \times Y$ und $(U \times Y) \cap (U' \times Y) = (U \cap U') \times Y = \emptyset$.
        
        Analog für $y \neq y'$. Damit ist $X \times Y$ hausdorffsch.

        \item \emph{Bemerkung.} Für $y \in Y$ ist $X \times \set{y} \subseteq X \times Y$ mit der Unterraumtopologie topologisch äquiavlent zu $X$ vermöge $(x,y) \mapsto x$.
        
        \emph{Beweis der Bemerkung.} Die Abildung 
        \begin{align*}
            X \times \set{y} \to& X \\ 
            (x,y) \mapsto& x
        \end{align*}
        ist die Einschränkung der Projektion $X \times Y \to X$ auf $X \times \set{y}$, also stetig. Die Umkehrfunktion 
        \begin{align*}
            X \to& X \times \set{y} \\
            x \mapsto& (x,y)
        \end{align*}
        ist stetig, da ihre Komponenten $\id_X: X \to X$ und $\text{konst}_y: X \to \set{y} \subseteq Y$ stetig sind.

        %%%%%%%%%%%%%%%%%%%%%%%%%%%%%%%%%%%%%%%%%%%%%%%%%%%%%%%%%%%%%%%%%%%%%%%%%%%%%%%%%%%%%%%%%%%%%%%%%%%%
        %%%%%%%%%%%%%%%%%%%%%%%%%%%%%%%%%%%%%%%%%% 2024 / 05 / 02 %%%%%%%%%%%%%%%%%%%%%%%%%%%%%%%%%%%%%%%%%%
        %%%%%%%%%%%%%%%%%%%%%%%%%%%%%%%%%%%%%%%%%%%%%%%%%%%%%%%%%%%%%%%%%%%%%%%%%%%%%%%%%%%%%%%%%%%%%%%%%%%%

        \breakline{Beginn Vorlesung 08 (02.05.2024)}

        Sei $\set{U_i}_{i \in I}$ eine offene Überdeckung von $X \times Y$.

        \textbf{Tubenlemma.}\index{Tubenlemma} Sei $y \in Y$. Dann gibt es eine offene Umgebung $U_y$ von $y$, sodass $X \times U_y$ von endlich vielen $U_i$ überdeckt wird.

        \emph{Beweis des Tubenlemmas.} Nach obiger Bemerkung ist $X \times \set{y}$ topologisch äquiavlent zu $X$, also ist $X \times \set{y}$ quasikompakt. Jedes $U_i$ ist von der Form 
        \begin{equation*}
            U_i := \bigcup_{j \in J_i} V_j \times W_j
        \end{equation*}
        mit $V_j \subseteq X$ und $W_j \subseteq Y$ offen. Dann ist $\set{V_j \times W_i}_{i \in I, j \in J_i}$ eine offene Überdeckung von $X \times Y$.

        Da $X \times \set{y}$ quasikompakt ist, gibt es $i_1, \ldots i_n \in I, j_1 \in J_{i_1}, \ldots, j_n \in J_{i_n}$ mit
        \begin{align*}
            X \times \set{y} \subseteq \bigcup_{k = 1}^n (V_{j_k} \times W_{j_k}).
        \end{align*}
        OBdA sei $y \in W_{j_k}$ für $k = 1,\ldots,n$. Setze
        \begin{equation*}
            U_y := \bigcap_{k = 1}^n W_{j_k}
        \end{equation*}
        offen mit $y \in U_y$. Sei $(x,y') \in X \times U_y$. Dann ist $(x,y)$ in einem $V_{j_k} \times W_{j_k}$. Dann ist
        \begin{equation}
            (x,y') \in V_{j_k} \times U_y \subseteq V_{j_k} \times W_{j_k} \subseteq U_{i_k}.
        \end{equation}
        Damit folgt 
        \begin{equation*}
            X \times U_y \subseteq \bigcup_{k = 1}^n U_{i_k}
        \end{equation*}
        und das Tubenlemma ist gezeigt.

        Sei nun für jedes $y \in Y$ die Menge $U_y \subseteq Y$ eine offene Umgebung wie im Tubenlemma. Dann ist $\set{U_y}_{y \in Y}$ eine offene Überdeckung von $Y$. Da $Y$ quasikompakt ist, gibt es eine endliche Teilüberdeckung 
        \begin{equation*}
            Y = \bigcup_{k = 1}^n U_{y_k}. 
        \end{equation*}
        Da jedes $X \times U_{y_k}$ von endlich vielen der $U_i$ überdeckt wird, wird auch 
        \begin{equation*}
            X \times Y = \bigcup_{k = 1}^n (X \times U_{y_k})
        \end{equation*}
        von endlich vielen der $U_i$ überdeckt und die Aussage ist gezeigt.
    \end{enumerate}
\end{proof}

Seien $X_1, \ldots, X_n$ topologische Räume. Wir versehen das Prodult $X_1 \times \cdots \times X_n$ induktiv mit der Produkttopologie vermöge $X_1 \times \cdots \times X_n = (X_1 \times \cdots \times X_{n-1}) \times X_n$.

\begin{lem}{Basis der Produkttopologie}{baseProductTopology}
    Seien $X, Y$ topologische Räume mit Basen $\cB_X$ und $\cB_Y$. Dann ist 
    \begin{equation*}
        \cB_{X \times Y} := \set{U \times V \mid U \in \cB_X, V \in \cB_Y}
    \end{equation*}
    eine Basis der Produkttopologie auf $X \times Y$.
\end{lem}
\begin{proof}
    Übungsaufgabe.
\end{proof}

\textbf{Korollar.} Die Blöcke $U_1 \times \cdots \times U_n$ mit $U_i \subseteq X_i$ offen bilden eine Basis der Produkttopologie auf $X_1 \times \cdots \times X_n$. 

Insbesondere ist die Produkttopologie unabhängig von der Klammerung.

\textbf{Korollar} (von Satz~\ref{thm:HausdorffQuasicompactProduct})\textbf{.} Seien $X_1, \ldots, X_n$ topologische Räume.
\begin{enumerate}[$(i)$]
    \item Wenn $X_i$ hausdorffsch ist für $i = 1,\ldots,n$, dann ist auch $X_1 \times \cdots \times X_n$ hausdorffsch.
    \item Wenn $X_i$ quasikompakt ist für $i = 1,\ldots,n$, dann ist auch $X_1 \times \cdots \times X_n$ quasikompakt.
\end{enumerate}

\begin{thm}{Heine-Borel}{heineBorel}
    Sei $A \subseteq \IR^n$. Dann ist $A$ genau dann kompakt, wenn $A$ abgeschlossen und beschränkt ist.
\end{thm}

\begin{proof}\leavevmode
    \begin{itemize}
        \item[\enquote{$\Rightarrow$}:] Die Aussage gilt für allgemeine metrische Räume $X$ und $A \subseteq X$.

        \emph{Beschränktheit:} Sei $x \in X$. Dann ist
        \begin{equation*}
            X = \bigcup_{n \in \IN^*} U(x,n).
        \end{equation*}
        Da $A$ kompakt ist, gibt es $n_1, \ldots, n_k \in \IN^*$ mit 
        \begin{equation*}
            A \subseteq \bigcup_{i = 1}^k U(x,n_i).
        \end{equation*}
        OBdA ist $n_1 \geq n_i$ für alle $i = 2,\ldots,k$. Dann ist $A \subseteq U(x,n_1)$, also ist $A$ beschränkt.

        \emph{Abgeschlossenheit:} Da $X$ ein metrischer Raum ist, ist $X$ ein Hausdorffraum. Da $A \subseteq X$ quasikompakt ist, ist nach $A$ Satz~\ref{thm:HausdorffClosedSubspace} auch abgeschlossen.

        \item[\enquote{$\Leftarrow$}:] Da $A$ beschränkt ist, gibt es ein $r > 0$ mit 
        \begin{equation*}
            A \subseteq \intervalcc{-r,r}^n subseteq \IR^n.
        \end{equation*}
        Da $\intervalcc{-r,r}$ topologisch äquiavlent zu $\intervalcc{0,1}$ ist, ist $\intervalcc{-r,r}$ kompakt (nach Satz~\ref{thm:compact01}). Mit Satz\ref{thm:HausdorffQuasicompactProduct} ist dann auch $\intervalcc{-r,r}^n$ kompakt, und nach Satz~\ref{thm:HausdorffClosedSubspace} in $\intervalcc{-r,r}^n$, also ist $A$ kompakt.
    \end{itemize}
\end{proof}

\begin{beispiel}[Endliche Dimension ist notwendig]
    Die Menge 
    \begin{equation*}
        \ell(\IN) := \set{(x_0, x_1, \ldots) \in \IR^\IN \mid \sum_{i = 0}^\infty x_i^2 < \infty}
    \end{equation*}
    ist ein $\IR$-Vektorraum bezüglich punktweiser Skalarmultiplikation und der Metrik 
    \begin{equation*}
        d\bigl((x_n)_{n \in \IN}, (y_n)_{n \in \IN}\bigr) := \sqrt{\sum_{n = 1}^\infty (x_n - y_n)^2}.
    \end{equation*}
    Die Teilmenge 
    \begin{equation*}
        \mathcal{K} := \set{(x_n)_{n \in \IN} \mid d\bigl((x_n)_{n \in \IN}, 0\bigr) \leq 1}
    \end{equation*}
    ist abgeschlossen (da $d$ stetig) und beschränkt, aber nicht kompakt:

    Die Folge $x_i = e_i$, also mit einer 1 am $i$-ten Eintrag und einer 0 an allen anderen Einträgen, hat keine konvergente Teilfolge.
\end{beispiel}

Sei $\set{X_i}_{i \in I}$ eine Familie topologischer Räume. Setze 
\begin{equation*}
    \prod_{i \in I} X_i := \set{(x_i)_{i \in I} \mid x_i \in X_i, i \in I}.
\end{equation*}

\begin{defi}{Produkttopologie (allgemein)}{productTopologyInf}
    Die \emph{Produkttopologie}\index{Produkttopologie}\index{Produkttopologie!über Familien} auf $\prod_{i \in I} X_i$ ist die Topologie erzeugt von 
    \begin{equation*}
        \cS := \set{U_j \times \prod_{i \in I \setminus \set{i}} X_i \mid j \in I, U_j \subseteq X_j \text{ offen}}.
    \end{equation*}
\end{defi}

\begin{bemerkung}
    $\cS$ ist eine Subbasis der Produkttopologie, aber im Allgemeinen keine Basis.

    Eine Basis ist gegeben durch 
    \begin{equation*}
        \cB := \set{\prod_{j \in J} U_j \times \prod_{i \in I \setminus J} X_i \mid \text{$J \subseteq I$ endlich, $U_j \subseteq X_j$ offen}}.
    \end{equation*}
\end{bemerkung}

\textbf{Korollar.} Falls $I$ endlich ist, stimmt die Defintion mit der Definition~\ref{def:productTopology} überein.

\textbf{Warnung.} Wenn $I$ unendlich ist, so sind Mengen der Form 
\begin{equation*}
    \prod_{i \in I} U_i, \quad \text{$U_i \subseteq X_i$ offen},
\end{equation*}
im Allgemeinen nicht offen in der Produkttopologie. 

\begin{thm}{Stetigkeit bezüglich Produkttopologie $\Leftrightarrow$ komponentenweise stetig}{continuousProductVsContinuousByComponent}
    Sei $(X_i)_{i \in I}$ eine Familie topologischer Räume. Sei $A$ ein topologischer Raum und seien $f_i: A \to X_i$ für $i \in I$ Funktionen. Definiere 
    \begin{align*}
        f: A \to& \prod_{i \in I} X_i \\
        a \mapsto& (f_i(a))_{i \in I}.
    \end{align*}
    Dann ist $f$ gnau dann stetig bezüglich der Produkttopologie auf $\prod_{i \in I} X_i$, wenn alle $f_i$ stetig sind.
\end{thm}

\begin{proof}\leavevmode
    \begin{itemize}
        \item[\enquote{$\Rightarrow$}:] Sei $j \in I, \pr_j: \prod_{i \in I} X_i \to X_j$ die Projektion. $\pr_j$ ist stetig: Sei $U_j \subseteq X_j$ offen. Dann ist 
        \begin{equation*}
            \pr_j^{-1}(U_j) := U_j \times \prod_{i \in I \setminus \set{j}} X_i
        \end{equation*}
        ein Element von $\cS$, insbesondere also offen.

        Also ist $f_j = \pr_j \circ f$ stetig als Verknüpfung stetiger Funktionen.

        \item[\enquote{$\Leftarrow$}:] Es genügt zu zeigen, dass $f^{-1}(Y)$ offen ist für alle $Y \in \cS$. Jedes $Y \in \cS$ ist von der Form
        \begin{equation*}
            Y = U_j \times \prod_{i \in I \setminus \set{j}} X_i
        \end{equation*}
        für $j \in I, U_j \subseteq X_j$ offen. Dann ist 
        \begin{equation*}
            f^{-1}(Y) = f_j^{-1}(U_j)
        \end{equation*}
        offen, da $f_j$ stetig ist.
    \end{itemize}
\end{proof}

\begin{beispiel}
    Seien $X_1, \ldots, X_n$ diskrete Räume. Dann ist auch $X_1 \times X_2 \times \cdots \times X_n$ diskret, da 
    \begin{equation*}
        \set{(x_1,\ldots,x_n)} = \underset{\substack{\text{offen in}\\ X_1}}{\underbrace{\set{x_1}}} \times \cdots \times \underset{\substack{\text{offen in}\\ X_n}}{\underbrace{\set{x_n}}}
    \end{equation*}
    in der Basis der Produkttopologie liegt.
\end{beispiel}

\begin{beispiel}
    Betrachte $\set{0,2}$ mit der diskreten Topologie. Dann ist 
    \begin{equation*}
        \prod_\IN \set{0,2} = \set{0,2}^\IN = \set{(a_i)_{i \in \IN} \mid a_i \in \set{0,2}}
    \end{equation*}
    topologisch äquivalent zu der \emph{Cantor-Menge}\index{Cantor-Mege}, einer nicht diskreten Teilmenge von $\IR$.
\end{beispiel}

\begin{definition}
    Definiere abgeschlossene Teilmengen $C_i \subseteq \intervalcc{0,1}$ induktiv. Setze $C_1 := \intervalcc{0,\frac{1}{3}} \cup \intervalcc{\frac{2}{3},1}$, und induktiv entfnen wir immer das mittlere Drittel (also $C_2 := \intervalcc{0,\frac{1}{9}} \cup \intervalcc{\frac{2}{9},\frac{1}{3}} \cup \intervalcc{\frac{2}{3},\frac{7}{9}}) \cup \intervalcc{\frac{8}{9}, 1}$.

    Die Menge $C_i$ besteht aus $2^i$ abgeschlossenen Intervallen der Länge $3^{-i}$. Die \emph{Cantor-Menge} $\cC$ ist definiert durch 
    \begin{equation*}
        \cC := \bigcup_{i \in I} C_i.
    \end{equation*}
\end{definition}
\emph{Alternativ:} Verwende die $3$-adische Darstellung einer rellen Zahl; jedes $x \in \intervalcc{0,1}$ lässt sich schreiben als 
\begin{equation*}
    x = \sum_{i = 1}^\infty \frac{a_i}{3^i}, \quad a_i \in \set{0,1,2}.
\end{equation*}
Es ist $\cC$ die Menge aller $x \in \intervalcc{0,1}$, wo alle Koeffizienten der 3-adischen Darstellung nur 0 oder 2 sind. 

\begin{bemerkung}
    Jedes $x$ hat höchstens eine Darstellung in dieser Form.
\end{bemerkung}
\begin{proof}
    Angenommen, $x = \sum_{i = 1}^\infty \frac{a_i}{3^i} = \sum_{i = 1}^\infty \frac{a_i'}{3^i}$ mit $a_i, a_i' \in \set{0,2}$ mit $a_i \neq a_i'$ für ein $i$. Sei $n := \min\set{i \in \IN \mid a_i \neq a_i'}$.  OBdA sei $a_n = 2, a_n' = 0$. Dann ist 

    \begin{equation*}
        0 = \sum_{i = 1}^\infty \frac{a_i}{3^i} - \sum_{i = 1}^\infty \frac{a_i'}{3^i} = \frac{2}{3^n} + \sum_{i = n+1}^\infty \frac{a_i - a_i'}{3^i}.
    \end{equation*}
    Es ist 
    \begin{equation*}
        \abs{\sum_{i = n+1}^\infty \frac{a_i - a_i'}{3^i}} \leq \sum_{i = n+1}^\infty \frac{\overset{\leq 2}{\overbrace{\abs{a_i - a_i'}}}}{3^i} \leq \sum_{i = n+1}^\infty \frac{2}{3^i} = \frac{2}{3^n} \sum_{i = 1}^\infty \frac{1}{3^i} = \frac{1}{3^n}
    \end{equation*}
    mit der geometrischen Reihe. Diese Zahl ist aber nicht darstellbar. $\lightning$
\end{proof}

Da $\cC$ abgeschlossen und beschränkt ist, ist $\cC$ kompakt nach Satz~\ref{thm:heineBorel}, aber nicht diskret, da $\set{0}$ nicht offen ist: Für jedes $\varepsilon > 0$ enthält $\cC \cap \intervalco{0,\varepsilon}$ unendlich viele Elemente von $\frac{2}{3^i}$.

Betrachte $f: \set{0,2}^\IN \to \cC, (a_i)_{i \in \IN} \mapsto \sum_{i = 1}^\infty \frac{a_i}{3^i}$.
\begin{thm}{$\set{0,2}^\IN \iso \cC$ (Cantormenge)}{cantorSet}
    $f$ (wie oben) ist eine topologische Äquivalenz.
\end{thm}
\begin{proof}
    $f$ ist nach Konstruktion bijektiv.

    $f^{-1}: \cC \to \set{0,2}^\IN$ ist stetig: Nach der Universellen Eigenschaft der Produkttopologie genügt es zu zeigen, dass jede Komponentenfunktion $(f^{-1})_i: \IC \to \set{0,2}$ stetig ist:

    Seien $x \in \left(f^{-1}\right)_i^{-1}(0), y \in \left(f^{-1}\right)_i^{-1}$. Da sich die 3-adische Darstellungen von $x$ und $y$ an $i$-ter Stelle um 2 unterscheiden, ist $\abs{y-x} \geq 3^{-i}$. Also ist $U(x, 3^{-i}) \cap \left(f^{-1}\right)_i^{-1}(2) = \emptyset = U(y, 3^{-i}) \cap \left(f^{-1}\right)_i^{-1}(0)$. Daraus folgt, dass $\left(f^{-1}\right)_i^{-1}(0)$ und $\left(f^{-1}\right)_i^{-1}(2)$ offen sind, also ist $(f^{-1})_i$ stetig, und damit auch $f^{-1}$.

    Da $\cC$ kompakt ist, ist $\set{0,2}^\IN$ hausdorffsch (\emph{allgemeiner:} Beliebige Produkte von Hausdorff-Räumen sind hausdorffsch mit gleichem Beweis wie im endlichen Fall (Satz~\ref{thm:HausdorffQuasicompactProduct})). Damit ist $f^{-1}: \cC \to \set{0,2}^\IN$ abgeschlossen und $f$ folglich stetig.
\end{proof}


%%%%%%%%%%%%%%%%%%%%%%%%%%%%%%%%%%%%%%%%%%%%%%%%%%%%%%%%%%%%%%%%%%%%%%%%%%%%%%%%%%%%%%%%%%%%%%%%%%%%
%%%%%%%%%%%%%%%%%%%%%%%%%%%%%%%%%%%%%%%%%% 2024 / 05 / 07 %%%%%%%%%%%%%%%%%%%%%%%%%%%%%%%%%%%%%%%%%%
%%%%%%%%%%%%%%%%%%%%%%%%%%%%%%%%%%%%%%%%%%%%%%%%%%%%%%%%%%%%%%%%%%%%%%%%%%%%%%%%%%%%%%%%%%%%%%%%%%%%

\breakline{Beginn Vorlesung 09 (07.05.2024)}

\begin{bemerkung}
    $\cC$ ist topologisch äquiavlent zu $\cC \times \cC$: Es ist 
    \begin{equation*}
        \cC \iso \prod_\IN \set{0,2} \iso \prod_{\substack{2\IN \\ \subseteq \IN}} \set{0,2} \times \prod_{\substack{2\IN - 1 \\ \subseteq \IN}} \set{0,2} \iso \prod_\IN \set{0,2} \times \prod_\IN \set{0,2} \iso \cC \times \cC.    
    \end{equation*}
\end{bemerkung}
\begin{bemerkung}[Satz: Tychonoff]
    Sei $(X_i)_{i \in I}$ eine Familie quasikompakter Räume. Dann ist auch $\prod_{i \in I} X_i$ quasikompakt.
\end{bemerkung}

% \textbf{Ziel:} 
\begin{thm}{Urysohnscher Metrisierungssatz}{urysohnMetric}
    Jeder normale Raum mit abzählbarer Basis der Topologie ist metrisierbar.
\end{thm}
\emph{Strategie:} 
\begin{enumerate}
    \item Der Raum 
    \begin{equation*}
        \prod_\IN \IR = \IR^\IN
    \end{equation*}
    als Produkttopologie ist metrisierbar,
    \item Jeder normale Raum mit abzählbarer Basis der Topologie ist topologisch äquiavlent zu einem Unterraum von $\IR^\IN$.
\end{enumerate}

Wir beweisen den Satz durch diese beiden Schritte.
\begin{enumerate}
    \item Definiere eine Metrk $d'$ auf $\IR$ durch 
    \begin{equation*}
        d'(x,y) := \min\set{\abs{x-y}, 1}.
    \end{equation*}
    \begin{lemm}
        $d'$ induziert die Standardtopologie af $\IR$.
    \end{lemm}
    Wir definieren die Metrik $D$ auf $\prod_\IN \IR$ durch 
    \begin{equation*}
        D\bigl((x_n)_{n \in \IN}, (y_n)_{n \in \IN}\bigr) := \sup\set{\frac{d'(x_n, y_n)}{n} \leq \frac{1}{n} \mid n \in \IN}.
    \end{equation*}
    Wir zeigen, dass $D$ ein Metrik ist. Dafür muss die Dreiecksungleichung gezeigt werden: Für jedes $n \in \IN$ ist 
    \begin{align*}
        D\bigl((x_n)_{n \in \IN}, (z_n)_{n \in \IN}\bigr) &=\frac{d'(x_n,z_n)}{n} \leq \frac{d'(x_n,y_n)}{n} + \frac{d'(y_n,z_n)}{n} \\
        &\leq D\bigl((x_n)_{n \in \IN}, (y_n)_{n \in \IN}\bigr) + D\bigl((y_n)_{n \in \IN}, (z_n)_{n \in \IN}\bigr)
    \end{align*}
    \begin{lem}{Produkttopologie auf $\IR^\IN$}{productTopolgyIRIN}
        $D$ induzert die Produkttopologie auf $\IR^\IN$.
    \end{lem}
    \begin{proof}
        Sei $U \subseteq \prod_\IN \IR$ offen bezüglich der metrischen Topologie, $x \in U$. Dann gibt es ein $\varepsilon > 0$, sodass $U_D(x,\varepsilon) \subseteq U$.

        Sei $N \in \IN$, sodass $\frac{1}{N} < \varepsilon$. Setze 
        \begin{equation*}
            V = \intervaloo{x_1 - \varepsilon, x_1 + \varepsilon} \times \intervaloo{x_2 - \varepsilon, x_2 + \varepsilon} \times \cdots \times \intervaloo{x_N - \varepsilon, x_N + \varepsilon} \times \IR \times \IR \times \cdots
        \end{equation*}
        $V$ ist offen in der Produkttopologie und es ist $x \in V$. 

        Sei $(y_n)_{n \in \IN}$. Dann ist 
        \begin{equation*}
            \frac{d'(x_n, y_n)}{n} < \begin{cases}
                \frac{\varepsilon}{n} & n \leq N, \\ 
                \frac{1}{N} < varepsilon & n > N.
            \end{cases}
        \end{equation*}
        Also ist $D\bigl((x_n)_{n \in \IN}, (y_n)_{n \in \IN}\bigr) < \varepsilon$. Damit ist $(y_n)_{n \in \IN} \in U_D(x, \varepsilon)$ und damit auch $x \in V \subseteq U_D(x, \varepsilon) \subseteq U$ und $U$ ist offen bezüglich der Produkttopologie.

        Umgekehrt sei $V \subseteq \prod_\IN \IR$ offen bezüglich der Produkttopologie, $x \in V$. Dann gibt es $V_1, \ldots, V_n \subseteq \IR$ offen mit 
        \begin{equation*}
            x \in V_1 \times \cdots \times V_n \times \IR \times \IR \times \cdots \quad \subseteq V.
        \end{equation*}
        Für $i = 1, \ldots, n$ sei $\varepsilon_i > 0$ mit $\intervaloo{x_i - \varepsilon_i, x_i + \varepsilon_i} \subseteq V_i$. Setze $\varepsilon := \min\set{\frac{\varepsilon_i}{i} \mid i = 1,\ldots,n}$. Dann ist 
        \begin{equation*}
            \frac{\abs{y_i - x_i}}{i} \leq D(x, y) < \frac{\varepsilon_i}{i} \quad \quad \text{für alle $i = 1,\ldots,n$}.
        \end{equation*}
        Damit folgt $\abs{y_i - x_i} < \varepsilon_i \Rightarrow y_i \in V_i$ für $i = 1, \ldots, n$, also folgt 
        \begin{equation*}
            y \in V_1 \times \cdots \times V_n \times \IR \times \IR \times \cdots \quad \subseteq V,
        \end{equation*}
        also $U_D(x, \varepsilon) \subseteq V$. Damit ist $V$ offen in der metrischen Topologie.
    \end{proof}
    \item Sei $X$ ein normaler Raum mit Basis $\cB := \set{B_1, B_2, \ldots}$.
    
    \begin{lem}{Abzählbare Familie stetiger Funktionen fast überall 0}{countableContinuous}
        Es gibt eine abzählbare Familie stetiger Funktionen $f_n: X \to \intervalcc{0,1}$, sodass für jedes $x \in X$ und für jede Umgebung $U$ von $x$ ein $n \in \IN$ existiert, sodass 
        \begin{align*}
            f_n(x) &> 0, \\
            f_n(y) &= 0 \quad \quad \text{für alle $y \notin U$}.
        \end{align*}
    \end{lem}
    \begin{proof}
        Für jedes Paar $n, m$ mit $\overline{B_n} \subseteq B_m$ wenden wir das Urysohnsche Lemma (Satz~\ref{thm:urysohn}) an, und wählen Funktionen
        \begin{align*}
            g_{n,m}&: X \to \intervalcc{0,1} \quad \text{mit} \\
            g_{n,m}&(\overline{B_n}) = \set{1} \quad \text{und} \\
            g_{n,m}&(X \setminus B_m) = \set{0}.
        \end{align*}
        Die Menge solcher Paare ist abzählbar und die Funktionen erfüllen das Gewünschte: Sei $x \in X, x \in U$ Umgebung von $x$. Damit gibt es ein $m \in \IN$ mit $x \in B_m \subseteq U$.

        Da $X$ normal ist, gibt es eine offene Menge $V$ mit $x \in V$ und $\overline{V} \subseteq B_m$. Sei $B_n$ so, dass $x \in B_n \subseteq V$. Dan ist 
        \begin{equation*}
            \overline{B_n} \subseteq \overline{V} \subseteq B_m
        \end{equation*}
        und es gilt $g_{n,m}(x) = 1$, da $x \in B_n \subseteq \overline{B_n}$, und für $y \notin U$ giklt $g_{n,m}(y) = 0$, da $y \notin B_m$.

        Seien $(f_n)_{n \in \IN}$ wie im Lemma~\ref{lem:countableContinuous}. Wir definieren 
        \begin{align*}
            F: X &\to \prod_\IN \IR \\
            x &\mapsto \bigl(f_n(x)\bigr)_{n \in \IN}.
        \end{align*}
        Nach der universellen Eigenschaft der Produkttopologie (Satz~\ref{thm:universalTraitProductTopo}) ist $F$ stetig.
        
    \end{proof}
    \emph{Behauptung.} $F$ ist eine Topologische Äquivalenz auf $F(X) := Z$.
    \begin{proof}
        $F$ ist injektiv: Seien $x \neq y \in X$. Da $X$ hausdorffsch ist, gibt es eine Umgebung $V$ von $x$ mit $y \notin U$. Nach Lemma~\ref{lem:countableContinuous} gibt es ein $n \in \IN$ mit $f_n(x) =  > 0$ und $f_n(X \setminus U) = \set{0}$. Damit ist $f_n(x) \neq f_n(y) \Rightarrow F(x) \neq F(y)$.

        Sei $U \subseteq X$ offen, also ist auch $F(U)$ offen in $Z = F(X)$. Sei $z \in F(U)$ mit eindeutigem Urbild $x$. Wir konstruieren $V \subseteq \prod_\IN \IR$ offen, sodass 
        \begin{equation*}
            z \in \underset{\text{offen in $F(X)$}}{\underbrace{V \cap F(X)}} \subseteq F(U).
        \end{equation*}
        Nach Lemma~\ref{lem:countableContinuous} gibt es $n \in \IN$ mit $f_n(x) > 0$ und $f_n(X \setminus U) = \set{0}$. Setze nun 
        \begin{equation*}
            \IR \times \IR \times \cdots \times \IR \times \underset{\text{$n$-te Komp.}}{\underbrace{\intervaloo{0,\infty}}} \times \IR \times \IR \times \cdots.
        \end{equation*}
        Die Menge $V$ ist offen in $\prod_\IN \IR$. Es ist 
        \begin{equation*}
            z_n = f_n(x) > 0 \Rightarrow z \in V \cap F(X). 
        \end{equation*}
        Sei $z' = F(x') \in V \cap F(X)$. Es gilt $x' \in U \Rightarrow z' = F(x') \in F(U)$, also folgt $V \cap F(X) \subseteq F(U)$. Damit ist $F(U)$ eine Umgebung von $z$, also insbesondere offen in $Z$.

        Damit ist Stetigkeit und somit auch topologische Äquivalenz gezeigt.
    \end{proof}
\end{enumerate}

\begin{bemerkung}
    \begin{enumerate}[$(i)$]
        \item Nicht jeder metrische Raum hat eine abzählbare Basis der Topologie.
        \begin{beispiel}
            $X$ überabzählbar, mit der diskreten Metrik ausgestattet.
        \end{beispiel}
        \item Jeder Punkt $x \in X$ eines metrischen Raums hat eine abzählbare Umgebungsbasis, d.h. es gibt Umgebungn $x \in U_n \subseteq X$ für $n \in \IN$ von $x$; sodass ede andere Umgebung $U$ eins der $U_n$ enthält.
        \begin{beispiel}
            Setze $U_n := U(x, \frac{1}{n})$.
        \end{beispiel}
        \item Es gibt normale Räume ohne abzählbare Umgebungsbasis.
        \begin{beispiel}
            Sei $X = \prod_I \intervalcc{0,1}$ mit $I$ überabzählbar. Diese Menge ist nach Tychonoff kompakt, aber nicht normal.
        \end{beispiel}
    \end{enumerate}
\end{bemerkung}


\chapter{Disjunkte Verenigungen, Verklebungen und (Weg-)Zusammenhang}

Sei $(X_i)_{i \in I}$ eine Familie von Mengen. Die \emph{disjunkte Vereinigung}\index{disjunkte Vereinigung} $\bigsqcup_{i \in I} X_i$ ist definiert als 
\begin{equation*}
    \bigsqcup_{i \in I} X_i := \set{(i, x) \mid i \in I, x \in X_i}.
\end{equation*}
Für jedes $j \in I$ ist 
\begin{align*}
    i_j: X_j &\inj \bigsqcup_{i \in I} X_i \\
    x &\mapsto (j,x)
\end{align*}
injektiv und induziert eine Bijektion 
\begin{equation*}
    X_j \iso \set{(j,x) \mid x \in X_j} \subseteq \bigsqcup_{i \in I} X_i.
\end{equation*}

Sei nun $(X_i)_{i \in I}$ eine Familie topologischer Räume.
\begin{defi}{Offene Teilmengen disjunkter Vereinigungen}{openSetDisjointUnion}
    Eine Teilmenge $U \subseteq \bigsqcup_{i \in I} X_i$ heißt genau dann \emph{offen}\index{disjunkte Vereinigung!offene Teilmenge}\index{offen}, wenn 
    \begin{equation*}
        i_j^{-1}(U) \subseteq X_j
    \end{equation*}
    für alle $j \in I$ offen ist.
\end{defi}
Außerdem ist die disjunkte Vereinigung topologischer Räume wieder ein topologischer Raum (Beweis: Übung).

\begin{bemerkung}
    Für jedes $j \in I$ ist 
    \begin{equation*}
        \set{(j,x) \mid x \in X_j} \quad \quad \text{offen in} \quad \bigsqcup_{i \in I} X_i
    \end{equation*}
    und $i_j: X_j \to \bigsqcup_{i \in I} X_i$ ist stetig und induziert eine topologische Äquivalenz 
    \begin{equation*}
        X_j \iso \set{(j,x) \mid x \in X_j},
    \end{equation*}
    aufgefasst mit der Unterraumtopologie.
\end{bemerkung}

\begin{thm}{Universelle Eigenschaft der disjunkten Vereinigung}{universalTraitDisjointTopo}\index{Universelle Eigenschaft!disjunkte Vereinigung}\index{Universelle Eigenschaft}
    Sei $(X_i)_{i \in  I}$ eine Familie topologischer Räume, $Y$ ein topologischer Raum und $f_j: X_j \to Y$ für alle $j \in I$. Definiere 
    \begin{align*}
        F: \bigsqcup_{i \in I} X_i & \to Y \\
        (i,x) &\mapsto f_i(x).
    \end{align*}
    Dann ist $F$ genau dann stetig, wenn alle $f_i$ stetig sind.
\end{thm}
\begin{proof}\leavevmode
    \begin{itemize}
        \item[\enquote{$\Rightarrow$}:] Es ist $f_j = F \circ i_j$, also stetig als Verknüpfung stetiger Funktionen.
        \item[\enquote{$\Leftarrow$}:] Sei $V \subseteq Y$ offen. Wir wollen zeige, dass $F^{-1}(V)$ offen in $\bigsqcup_{i \in I} X_i$ offen ist. 

        Es ist $i_j^{-1}(F^{-1}(V)) = (F \circ i_j)^{-1}(V) = f_j^{-1}(V)$. Dies ist offen in $X_j$, da $f_j$ stetig ist, also ist auch $F^{-1}(V) \subseteq \bigsqcup_{i \in I} X_i$ offen, also ist $F$ stetig.
    \end{itemize}
\end{proof}

Sei nun $X$ ein topologischer Raum mit Unterräumen $X_1, X_2 \subseteq X$, wobei $X_1 \cup X_2 = X$.

\textbf{Frage.} Lässt sich die Topologie auf $X$ aus den Topologien auf $X_1, X_2$ und $X_\cap := X_1 \cap X_2$ rekonstriueren?

\begin{beispiel}
    Falls $X_1 \cap X_2 = \emptyset$ erhalten wir eine Abbildung $X_1 \sqcup X_2 \to X$, diese ist stetig nach Satz~\ref{thm:universalTraitDisjointTopo} und bijektiv. Dies ist genau dann eine topologische Äquivalenz, wen $X_1, X_2$ offen in $X$ sind.

    Sei beispielsweise $X = \intervalcc{0,1}, X_1 = \intervalcc{0,\frac{1}{2}}, X_2 = \intervaloc{\frac{1}{2},1}$. Dann ist 
    \begin{equation*}
        \intervalcc{0,\frac{1}{2}}~\dot{\cup}~\intervaloc{\frac{1}{2},1} \to \intervalcc{0,1}
    \end{equation*}
    keine topologische Äquivalenz.
\end{beispiel}

\emph{Allgemeiner:} Seien $X_0, X_1, X_2$ topologische Räume mit stetigen Abbildungen $f_1: X_0 \to X_1$ und $f_2: X_0 \to X_2$.

\emph{Verklebung.} Definiere $X_1 \cup_{X_0} X_2$ als Quotient 
\begin{equation*}
    \nicefrac{X_1 \sqcup X_2}{\sim}, \quad \quad \text{mit $\sim$ der von $f_1(x) \sim f_2(x)$ für alle $x \in X_0$ erzeugten Äquivalenzrelation.}
\end{equation*}
\emph{Vorsicht:} $X_1 \cup_{X_0} X_2$ hängt auch von $f_1$ und $f_2$ ab, obwohl sie nicht mehr in der Notation auftauchen.




%%%%%%%%%%%%%%%%%%%%%%%%%%%%%%%%%%%%%%%%%%%%%%%%%%%%%%%%%%%%%%%%%%%%%%%%%%%%%%%%%%%%%%%%%%%%%%%%%%%%
%%%%%%%%%%%%%%%%%%%%%%%%%%%%%%%%%%%%%%%%%% 2024 / 05 / 14 %%%%%%%%%%%%%%%%%%%%%%%%%%%%%%%%%%%%%%%%%%
%%%%%%%%%%%%%%%%%%%%%%%%%%%%%%%%%%%%%%%%%%%%%%%%%%%%%%%%%%%%%%%%%%%%%%%%%%%%%%%%%%%%%%%%%%%%%%%%%%%%

\breakline{Beginn Vorlesung 10 (14.05.2024)}

\begin{beispiel}
    Seien $X_1, X_2 = \IR$ und $X_0 = \IR \setminus \set{0}$ mit $f_1, f_2$ der Inklusion. Dann ist 
    \begin{equation*}
        \IR \cup_{\IR \setminus \set{0}} \IR
    \end{equation*}
    die \emph{reelle Gerade mit doppeltem Ursprung}.

    Dieser Raum ist nicht hausdorffsch, aber jeder Punkt hat eine Umgebung topologisch äquiavlent zu $\IR$.
\end{beispiel}

Wir betrachten folgende Diagramme:
\begin{equation*}
    \begin{tikzcd}
        X_1 \arrow[r, "i_1"] \arrow[rr, bend right, "j_1"] & X_1 \cup X_2 \arrow[r, "q"] & X_1 \cup_{X_0} X_2
    \end{tikzcd}
\end{equation*}

und 

\begin{equation*}
    \begin{tikzcd}
        X_2 \arrow[r, "i_2"] \arrow[rr, bend right, "j_2"] & X_1 \cup X_2 \arrow[r, "q"] & X_1 \cup_{X_0} X_2.
    \end{tikzcd}
\end{equation*}

\begin{lem}{Verklebung unter injektiven Funktionen}{attachingInjectve}
    \begin{enumerate}[$(i)$]
        \item Ist $f_1: X_0 \to X_1$ injektiv, so ist auch $j_2$ injektiv,
        \item Ist $f_2: X_0 \to X_2$ injektiv, so ist auch $j_1$ injektiv.
    \end{enumerate}
\end{lem}
\begin{proof}
    Aus Symmetriegründen zeigen wir nur $(i)$. Seien $x,y \in X_2$ mit $j_2(x) = j_2(y)$. Dann gibt es Punkte $x_1, \ldots, x_n \in X_0$, sodass 
    \begin{equation*}
        x = f_2(x_1), \quad f_1(x_1) = f_1(x_2), \quad f_2(x_2) = f_2(x_3), \quad \ldots, \quad f_2(x_n) = y.
    \end{equation*}
    Da $f_1$ injektiv ist, ist $x_1 = x_2, x_3 = x_4, x_5 = x_6$ usw., und damit folgt
    \begin{equation*}
        x = f_2(x_1) = f_2(x_2) = f_2(x_3) = \cdots = f_2(x_n) = y,
    \end{equation*}
    also ist $j_2$ injektiv.
\end{proof}

Sei nun $X$ ein topologischer Raum, $X_1, X_2 \subseteq X$ Unterräume mit $X_1 \cup X_2 = X$ und $X_0 := X_1 \cap X_2$. Seien $f_1: X_0 \to X_1$ und $f_2: X_0 \to X_2$ die Inklusionen. Betrachte 
\begin{align*}
    f': X_1 \sqcup X_2 &\surj X \\
    (1,x) &\mapsto x, \\
    (2,x) &\mapsto x.
\end{align*}
$f'$ faktorisiert über $f: X_1 \cup_{X_0} X_2 \to X$ und $f$ ist stetig wegen der universellen Eigenschaft sowie bijektiv (nach Lemma~\ref{lem:attachingInjectve} und Satz~\ref{thm:universalTraitDisjointTopo}).

\begin{thm}{Verklebung Topologische Äquivalenz}{attachingTopEquiv}
    $f$ ist eine topologische Äquivalenz, falls 
    \begin{enumerate}[$(i)$]
        \item $X_1$ und $X_2$ offen sind, oder
        \item $X_1$ und $X_2$ abgeschlossen sind.
    \end{enumerate}
\end{thm}
\begin{proof}
    Wir zeigen aus Symmetriegrüden nur $(ii)$. Es genügt zu zeigen, dass $f$ eine abgeschlossene Abbildung ist:

    Sei $A \subseteq X_1 \cup_{X_0} X_2$ abgeschlossen. Dann sind $j_1^{-1}(A) \subseteq X_1$ und $j_2^{-1}(A) \subseteq X_2$ abgeschlossen. Dann ist 
    \begin{equation*}
        f(A) = \underset{\subseteq X_1 \subseteq X}{\underbrace{j_1^{-1}(A)}} \cup \underset{\subseteq X_2 \subseteq X}{\underbrace{j_2^{-1}(A)}}.
    \end{equation*}
    Da $X_1$ abgeschlossen in $X$ ist, ist $j_1^{-1}(A)$ abgeschlossen in $X$, analog für $X_2$. Damit ist $f(A)$ ebenfalls abgeschlossen in $X$ als Vereinigung zweier in $X$ abgeschlossener Mengen.
\end{proof}

\begin{beispiel}
    Seu $X = \IS^n$. Setze $X_1 = \set{x \in \IS^n \mid x_{n+1} \geq 0}$ als \emph{obere Hemisphäre}, und $X_2 = \set{x \in \IS^n \mid x_{n+1} \leq 0}$ als \emph{untere Hemisphäre}. Dann ist 
    \begin{equation*}
        X_0 = X_1 \cap X_2 = \set{x \in \IS^n \mid x_{n+1} = 0} = \IS^{n-1} \times \set{0}.
    \end{equation*} 
    Nach Satz~\ref{thm:attachingTopEquiv} ist $\IS^n \iso X_1 \cup_{X_0} X_2$. Es sind $X_1$ und $X_2$ topologisch äquiavlent zu $D^n = \set{x \in \IR^n \mid \norm{x} \leq 1}$. Betrachte die Projektion
    \begin{align*}
        \varphi: X_1 &\bij D^n,  \\
        (x_1,\ldots,x_{n+1}) &\mapsto (x_1, \ldots, x_n)
    \end{align*}
    mit Umkehrfunktion 
    \begin{align*}
        \psi: D^n &\to X_1, \\
        (x_1, \ldots, x_n) &\mapsto \left(x_1, \ldots, x_n, \sqrt{1 - x_1^2 - \cdots - x_n^2}\right),
    \end{align*}
    analog für $X_2$.

    Also ist $\IS^n$ topologisch äquivalent zu $D^n \cup_{\IS^{n-1}} D^n$.
\end{beispiel}

\begin{defi}{Zusammenhang}{connected}
    Ein topologischer Raum $X$ heißt \emph{zusammenhängend}\index{Zusammenhang}\index{zusammenhängend}, wenn es keine Teilenge $A \subseteq X$ mit $A \neq \emptyset, X$ gibt, die offen und abgeschlossen sind.
\end{defi}
\begin{bemerkung}
    $X$ ist genau dann  zusammenhängend, wenn $X$ nicht topologisch äquivalent zu $X_1 \sqcup X_2$ mit $X_1, X_2 \neq \emptyset$ ist.
\end{bemerkung}
\begin{beispiel}\leavevmode
    \begin{enumerate}[$(i)$]
        \item $\IR \setminus \set{0} = \intervaloo{-\infty,0} \cup \intervaloo{0,\infty}$ ist nicht zusammenhängend, da die beiden Intervalle offen und disjunkt sind.
        \item $\IQ = \left(\IQ \cap \intervaloo{-\infty, \sqrt{2}}\right) \cup \left(\IQ \cap \intervalcc{\sqrt{2}, \infty}\right)$ ist nicht zusammenhängend, da beide Mengen offen und disjunkt sind.
    \end{enumerate}
\end{beispiel}

\begin{thm}{$\intervalcc{0,1}$ zusammenhängend}{unitIntervalConnected}
    Das Intervall $\intervalcc{0,1}$ ist zusammenhängend.
\end{thm}

\begin{proof}
    Angenommen, $\intervalcc{0,1} = A \cup B$ mit $A, B$ offen und disjunkt. Sei OBdA $0 \in A$ und $B \neq \emptyset$. Sei $t := \inf(B)$. Da $B$ abgeschlossen ist, ist $t \in B$. Damit ist $\intervalco{0,t} \subseteq A \quad \lightning$ zu $B$ offen.
\end{proof}

\begin{defi}{Wegzusammenhang}{pathconnected}
    Ein topologischer Raum $X$ heißt \emph{wegzusammenhängend}\index{Wegzusammenhang}\index{wegzusammenhängend}\index{Zusammenhang!Weg-}\index{zusammenhängend!weg-}, falls für je zwei Punkte $x, y \in X$ ein \enquote{Weg von $x$ nach $y$ existiert}, also falls eine stetige Funktion $w: \intervalcc{0,1} \to X$ existiert, sodass $w(0) = x$ und $w(1) = y$.
\end{defi}

\begin{lem}{(Weg-)Zusammenhang unter Surjektionen}{connectedUnderSurjection}
    Sei $f: X \surj Y$ stetig und surjektiv.
    \begin{enumerate}[$(i)$]
        \item Wenn $X$ zusammenhängend ist, ist auch $Y$ zusammenhängend.
        \item Wenn $X$ wegzusammenhängend ist, ist auch $Y$ wegzusammenhängend.
    \end{enumerate}
\end{lem}
\begin{proof}
    \begin{enumerate}[$(i)$:]
        \item Angenommen, $Y = U_1 \cup U_2$ mit $U_1, U_2 \subseteq Y$ offen, disjunkt und nichtleer. Dann ist 
        \begin{equation*}
            X = f^{-1}(Y) = \underset{\text{offen ($f$ stet.)}}{\underbrace{f{-1}(U_1)}} \cup \underset{\text{offen ($f$ stet.)}}{\underbrace{f^{-1}(U_2)}}.
        \end{equation*}
        Weiter ist 
        \begin{equation*}
            f^{-1}(U_1) \cap f^{-1}(U_2) = f^{-1}(U_1 \cap U_2) = \emptyset.
        \end{equation*}
        Da $f$ surejektiv, ist $f^{-1}(U_i) \neq \emptyset$ für $i = 1,2$. $\quad \lightning$ zu Zusammenhang von $X$.
        \item Seien $y_1, y_2 \in Y$. Da $f$ surjektiv, gibt es $x_1, x_2 \in X$ mit $f(x_1) = y_1, f(x_2) = y_2$.

        Da $X$ wegzusammenhängend ist, gibt es $w: \intervalcc{0,1} \to X$ stetig mit $w(0) = x_1, w(1) = x_2$. Dann ist $f \circ w: \intervalcc{0,1} \to Y$ stetig als Verknüpfung stetiger Funktionen und es gilt 
        \begin{align*}
            (f \circ w)(0) = f(w(0)) &= f(x_1) = y_1, \\
            (f \circ w)(1) = f(w(1)) &= f(x_2) = y_2.
        \end{align*}
        Also ist $f \circ w$ ein Weg von $y_1$ nach $y_2$.
    \end{enumerate}
\end{proof}

\begin{thm}{Wegzusammenhängend $\Rightarrow$ zusammenhängend}{pathconnectedImpConnected}
    Jeder wegzusammenhängende Raum ist auch zusammenhängend.
\end{thm}
\begin{proof}
    Sei $X$ wegzusammenhängend. Angenommen, $X = U_1 \cup U_2$ mit $U_1, U_2 \subseteq X$ offen, disjunkt und nichtleer. Sei $x_1 \in U_1, x_2 \in U_2$ und sei $w: \intervalcc{0,1} \to X$ stetig mit $w(0) = x_1, w(1) = x_2$. 

    Dann ist $w(\intervalcc{0,1})$ nach Lemma~\ref{lem:connectedUnderSurjection} zusammenhängend. 

    Andererseits ist 
    \begin{equation*}
        w(\intervalcc{0,1}) = \bigl(U_1 \cap w(\intervalcc{0,1})\bigr) \cup \bigl(U_2 \cap w(\intervalcc{0,1})\bigr),
    \end{equation*}
    wobei beide Mengen disjunkt und offen in $w(\intervalcc{0,1})$ sind. 

    Es ist $x_1 \in U_1 \cap w(\intervalcc{0,1})$ und $x_2 \in U_2 \cap w(\intervalcc{0,1})$, also sind beide nichtleer $\quad \lightning$ zum Zusammenhang von $w(\intervalcc{0,1})$.
\end{proof}

\begin{beispiel}
    $\IS^n$ ist wegzusammenhängend für $n \geq 1$, also auch zusammenhängend: Betrachte 
    \begin{align*}
        \IR^{n+1} \setminus \set{0} &\surj \IS^n \\
        x &\mapsto \frac{x}{\norm{x}}
    \end{align*}
    stetig und surjektiv. Im Beispiel $(ii)$ nach Definition~\ref{def:piZero} haben wir gesehen, dass $\IR^n \setminus \set{0}$ wegzusammenhängend ist, mit Lemma~\ref{lem:connectedUnderSurjection} ist $\IS^n$ wegzusammenhängend. 

    Außerdem ist $\IR^n$ wegzusammenhängend für alle $n$.
\end{beispiel}

\textbf{Erinnerung.} Sei $X$ ein topologischer Raum. Dann ist $\pi_0(X)$\index{$\pi_0(X)$ (Menge Weg-ZHK)} die Menge der Wegeäquivalenzklassen in $X$ (Definition~\ref{def:piZero}). Es gilt 
\begin{itemize}
    \item $\abs{\pi_0(\IR^n \setminus \set{0})} = 1$,
    \item $\abs{\pi_0(\IR \setminus \set{0})} = 2$;
    für $n \geq 2$. Damit ist $\IR$ nicht topologisch äquivalent zu $\IR^n$ für $n \geq 2$. (Siehe Satz~\ref{thm:RRnNotEquiv}).
    \item $\abs{\pi_0(\IS^1 \setminus \set{x,y})} = 2$ für je zwei $x,y \in \IS^1$, aber 
    \item $\abs{\pi_0(\IS^n \setminus \set{x,y})} = 1$ für je zwei $x,y \in \IS^n$ für $n \geq 2$.
\end{itemize}

\textbf{Korollar.} $\IS^1$ ist nicht topologisch äquiavlent zu $\IS^n$ für $n \geq 2$.



%%%%%%%%%%%%%%%%%%%%%%%%%%%%%%%%%%%%%%%%%%%%%%%%%%%%%%%%%%%%%%%%%%%%%%%%%%%%%%%%%%%%%%%%%%%%%%%%%%%%
%%%%%%%%%%%%%%%%%%%%%%%%%%%%%%%%%%%%%%%%%% 2024 / 05 / 16 %%%%%%%%%%%%%%%%%%%%%%%%%%%%%%%%%%%%%%%%%%
%%%%%%%%%%%%%%%%%%%%%%%%%%%%%%%%%%%%%%%%%%%%%%%%%%%%%%%%%%%%%%%%%%%%%%%%%%%%%%%%%%%%%%%%%%%%%%%%%%%%

\breakline{Beginn Vorlesung 11 (16.05.2024)}

\chapter*{Zusammenfassung \& Ausblick}

\textbf{Frage.} Wie lassen sich topologische Räume voneinander unterscheiden?

\emph{Bisher:}
\begin{description}
    \item[Eigenschaften] \begin{itemize}
        \item Quasikompaktheit
        \item Hausdorffeigenschaft
        \item Normalität
        \item (Weg-)Zusammenhang
    \end{itemize}
    \item[Invarianten] \begin{itemize}
        \item $\pi_0(X)$: Die Menge der Wegeäquivalenzklassen. Wenn $X$ topologisch äquiavlent zu $Y$ ist, haben $\pi_0(X)$ und $\pi_0(Y)$ gleich viele Elemente.
        \item \emph{Neu:} $\pi_1(X,x)$ ist $\nicefrac{\set{\text{Schleifen in $X$ an $x$}}}{\text{Deformationen bzw. Homotopien}}$, die \emph{Fundamentalgruppe}.
        
        Dabei wird die Multiplikation als Hintereinanderlaufen von Schleifen betrachtet. \emph{Jede Gruppe} kommt als Fundamentalgruppe vor! Ist $f: X \to Y$ eine topologische Äquivalenz, so sind $\pi_1(X,x)$ und $\pi_1(Y,f(x))$ isomorph. Es ist sogar \emph{Homotopieäquivalenz} genug.
        \begin{beispiel}\leavevmode
            \begin{itemize}
                \item $\IR$ ist homotopieäquivalent zu einem Zylinder,
                \item $\pi_1(\IS^1, x) \iso \IZ$ (\emph{Windungszahl}),
                \item $\pi_1(\IS^1 \times \IS^1, x) \iso \IZ \times \IZ$,
                \item $\pi_1(\IS^2, x) \iso \set{1}$,
                \item $\pi_1(\IR P^2, x) \iso \nicefrac{\IZ}{2\IZ}$.
            \end{itemize}
        \end{beispiel}
        \textbf{Korollar.} $\IS^1, \IS^1 \times \IS^1, \IS^2, \IR P^2$ sind paarweise nicht topologisch äquivalent.
    \end{itemize}
    \item[Hilfsmittel] zur Berechnung von $\pi_1(X,x)$: \emph{Überlagerungstheorie}. Beispielsweise $\exp: \IR \to \IS^1$.
\end{description}

\part{Algebraische Topologie}
\chapter{Die Fundamentalgruppe}

Sei $X$ ein topologischer Raum, $x \in X$ ein Punkt (\emph{Basispunkt}).

\begin{defi}{Schleife}{loop}
    Eine \emph{Schleife}\index{Schleife} in $X$ an $x$ ist eine stetige Funktion $w: \intervalcc{0,1} \to X$ mit $w(0) = w(1) = x$.
\end{defi}

\begin{defi}{Homotopie}{homotopy}
    Zwei Schleifen $w_1, w_2$ heißen \emph{homotop}\index{homotop}\index{Homotopie}, falls eine stetige Abbildung (\emph{Homotopie})
    \begin{equation*}
        H: \intervalcc{0,1} \times \intervalcc{0,1} \to X
    \end{equation*}
    gibt, sodass 
    \begin{enumerate}[$(i)$]
        \item $H(t, 0) = w_1(t)$ für alle $t \in \intervalcc{0,1}$,
        \item $H(t, 1) = w_2(t)$ für alle $t \in \intervalcc{0,1}$,
        \item $H(0, s) = H(1, s) = x$ für alle $s \in \intervalcc{0,1}$.
    \end{enumerate}
\end{defi}

\begin{thm}{Homotopie ist Äquivalenzrelation}{homotopyEquivalence}
    Homotopie ist eine Äquivalenzrelation auf der Menge der Schleifen in $X$ an $x$.
\end{thm}
\begin{proof}
    Wir müssen dafür die drei Eigenschaften prüfen:
    \begin{itemize}
        \item[\emph{$(i)$: Reflexivität.}] Sei $w$ eine Schleife an $x$. Dann ist $H(t,s) = w(t)$ eine Homotopie von $w$ zu $w$. 
        
        \emph{Formal:} $H$ ist die Verknüpfung 
        \begin{equation*}
            \begin{tikzcd}
                \intervalcc{0,1} \times \intervalcc{0,1} \arrow[r, "\pr_1"] \arrow[rr, bend right, "H2"] & \intervalcc{0,1} \arrow[r, "w"] & X.
            \end{tikzcd}
        \end{equation*}
        und damit stetig.
        \item[\emph{$(ii)$: Symmetrie.}] Sei $w_1$ homotop zu $w_2$ vermöge eine Homotopie $H$. Dann ist 
        \begin{equation*}
            H'(t,s) := H(t, 1-s)
        \end{equation*}
        also die Verknüpfung
        \begin{equation*}
            \begin{tikzcd}
                \intervalcc{0,1} \times \intervalcc{0,1} \arrow[r] & \intervalcc{0,1} \times \intervalcc{0,1} \arrow[r, "H"] & X \\
                (t,s) \arrow[r] & (t, 1-s)
            \end{tikzcd}
        \end{equation*}
        eine Homotopie von $w_2$ zu $w_1$.
        \item[\emph{$(iii)$: Transitivität.}] Seien $w_1$ homotop zu $w_2$ vermöge $H_1$ und $w_2$ homotop zu $w_3$ vermöge $H_2$. Dann ist 
        \begin{equation*}
            H(t,s) := \begin{cases}
                H_1(t, 2s) & s \in \intervalcc{0,\frac{1}{2}} \\
                H_2\left(t, 2\left(s - \frac{1}{2}\right)\right) & s \in \intervalcc{\frac{1}{2}, 1}
            \end{cases}
        \end{equation*}
        wohldefiniert, da $H\left(t, \frac{1}{2}\right) = H_1(t, 1) = w_2(t) = H_2(t, 0) = H_2\left(t, 2 \cdot \left(\frac{1}{2} - \frac{1}{2}\right)\right)$.

        Außerdem ist $H$ stetig, da $H|_{\intervalcc{0,1} \times \intervalcc{0,\frac{1}{2}}}$ und $H|_{\intervalcc{0,1} \times \intervalcc{\frac{1}{2}, 1}}$ stetig sind und beide Rechtecke abgeschlossen in $\intervalcc{0,1} \times \intervalcc{0,1}$ sind. Also ist $H$ eine Homotopie von $w_1$ zu $w_3$.
    \end{itemize}
    Damit ist die Aussage gezeigt.
\end{proof}

Wir schreiben $\pi_1(X, x)$ für 
\begin{equation*}
    \pi_1(X,x) := \nicefrac{\set{\text{Schleifen in $X$ an $x$}}}{\text{Homotopie}},
\end{equation*}
die \emph{Fundamentalgruppe von $X$ an $x$}.\index{Fundamentalgruppe}\index{$\pi_1(X, x)$ (Fundamentalgruppe)}

\emph{Ziel:} Gruppenstruktur auf $\pi_1(X,x)$. Seien $w, w'$ Schleifen in $X$ an $x$. Wir definieren eine neue Schleife 
\begin{equation*}
    (w * w')(t) := \begin{cases}
        w(2t) & 0 \leq t \leq \frac{1}{2}, \\
        w(2t - 1) & \frac{1}{2} \leq t \leq 1.
    \end{cases}
\end{equation*}
Diese realisiert das \enquote{Hintereinanderlaufen von Schleifen}.

\begin{thm}{Hintereinanderlaufen von Schleifen ist wohldefiniert}{loopToLoopWelldef}
    Das Hintereinanderlaufen von Schleifen ist wohldefiniert auf Homotopieklassen, und assoziativ nach Übergang zu Homotopieklassen, d.h. 
    \begin{align*}
        \pi_1(X, x) \times \pi_1(X, x) & \longrightarrow \pi_1(X,x) \\
        ([w], [w']) &\longmapsto [w * w']
    \end{align*}
    ist wohldefiniert und assoziativ.
\end{thm}

\begin{proof}
    Wir zeigen zuerst die \emph{Wohldefiniertheit}. Sei $H: \intervalcc{0,1} \times \intervalcc{0,1} \to X$ eine Homotopie von $w_0$ zu $w_1$, und $H': \intervalcc{0,1} \to \intervalcc{0,1} \to X$ eine Homotopie von $w_0'$ zu $w_1'$. Wir definieren 
    \begin{align*}
        H * H' : \intervalcc{0,1} \times \intervalcc{0,1} &\longrightarrow X \\
        (t,s) &\longmapsto \begin{cases}
            H(2t, s) & 0 \leq t \leq \frac{1}{2}, \\
            H'(2t - 1, s) & \frac{1}{2} \leq t \leq 1.
        \end{cases}
    \end{align*}
    Dies ist eine Homotopie von $w_0 * w_0'$ zu $w_1 * w_1'$. Also ist die Verknüpfung $*$ wohldefiniert auf Homotopieklassen.

    Zur \emph{Assoziativität:} Seien $w, w', w''$ Schleifen an $x$. Es ist 
    \begin{align*}
        H(t,s) = \begin{cases}
            w\left(\frac{4t}{1+2}\right) & 0 \leq t \leq \frac{1+s}{4} \\
            w'(4t-1-s) & \frac{1+s}{4}\leq t \leq \frac{2+s}{4} \\
            w''\left(\frac{4t-2-s}{2-s}\right) & \frac{2+s}{4} \leq t \leq 1
        \end{cases}
    \end{align*} 
    eine Homotopie von $(w * w') * w''$ zu $w * (w' * w'')$.
\end{proof}

\begin{thm}{Konstante Schleife ist neutrales Element, Inverses existiert}{fundamentalGroupIsActuallyGroup}
    Die Klasse der konstanten Schleife an $x$ ist ein neutrales Element für die Verknüpfung $*$ auf $\pi_1(X,x)$ und jedes Element hat ein Inverses (\enquote{Schleifen rückwärts laufen}).
\end{thm}
\begin{proof}
    Sei $k: \intervalcc{0,1} \to X$ mit $k(t) = x$ für alle $t \in \intervalcc{0,1}$ die \emph{konstante Schleife}. Sei $w: \intervalcc{0,1} \to X$ eine Schleife an $x$. Dann ist 
    \begin{equation*}
        H(t, s) := \begin{cases}
            w\left(\frac{2t}{1+s}\right) & 0 \leq t \leq \frac{1+s}{2} \\
            x & \frac{1+s}{2} \leq t \leq 1
        \end{cases}
    \end{equation*}
    eine Homotopie von $w * k$ zu $w$. Analog ist $k * w$ zu $w$. Also ist $[k]$ ein neutrales Element.

    \emph{Für das Inverse:} Wir definieren eine neue Schleife $\overline{w}$ durch $\overline{w}(t) := w(1-t)$. Dann ist 
    \begin{equation*}
        H(t, s) := \begin{cases}
            w(2t) & 0 \leq t \leq \frac{1-s}{2} \\
            w(1 - s) & \frac{1-s}{2} \leq t \leq \frac{1+s}{2} \\
            \overline{w}(2(1-t)) & \frac{1+s}{2} \leq t \leq 1
        \end{cases}
    \end{equation*}
    eine Homotopie von $w * \overline{w}$ zu $k$, also ist 
    \begin{equation*}
        [w] * [\overline{w}] = [w * \overline{w}] = [k].
    \end{equation*}
    Analog ist $[\overline{w}] * [w] = [k]$.
\end{proof}

\textbf{Korollar.} $(\pi_1(X,x), *)$ ist eine Gruppe.

\begin{thm}{$x, x'$ wegeäquivalent $\Rightarrow \pi_1(X,x) \iso \pi_1(X,x')$}{fundamentalGroupIsomorphism}
    Seien $x, x'$ zwei Punkte aus $X$, die durch einen Weg verbunden werden können. Dann sind $\pi_1(X,x)$ und $\pi_1(X, x')$ isomorphe Gruppen.
\end{thm}

\textbf{Korollar.} Fundamentalgruppen in verschiedenen Wegekomponenten haben im Allgemeinen nichts miteinander zu tun.

\begin{proof}
    Sei $v$ ein Weg von $v(0) = x$ zu $v(1) = x'$. Wir definieren eine Abbildung 
    \begin{align*}
        V_*: \pi_1(X, x') &\longrightarrow \pi_1(X, x) \\
        [w] &\longmapsto [v * w * \overline{v}],
    \end{align*}
    wobei $v * w * \overline{v}$ definiert ist durch 
    \begin{equation*}
        (v * w * \overline{v}) (t) = \begin{cases}
            v(3t) & 0 \leq t \leq \frac{1}{3} \\
            w(3t-1) & \frac{1}{3} \leq t \leq \frac{2}{3} \\
            v(3(1-t)) & \frac{2}{3} \leq t \leq 1.
        \end{cases}
    \end{equation*}
    Zu prüfen sind 
    \begin{description}
        \item[Wohldefiniertheit:] Es ist $w \hmtp w' \Rightarrow v * w * \overline{v} \hmtp v * w' * \overline{v}$.
        \item[$V_*$ ist Gruppenhomomorphismus:] Zu zeigen ist, dass $v * (w * w') * \overline{v} \hmtp (v * w * \overline{v}) * (v * w' * \overline{v})$: Es gilt 
        \begin{align*}
            (v * w * \overline{v}) * (v * w' * \overline{v}) \hmtp (v * w) * \underset{\hmtp k}{\underbrace{\overline{v} * v}} * (w' * \overline{v}) \hmtp (v * w) * (w' * \overline{v}) hmtp v * (w * w') * \overline{v}.
        \end{align*} 
        \item{Bijektivität von $V_*$} Sei $\overline{v}: \intervalcc{0,1} \to X$ der inverse Weg $\overline{v}(t) = v(1-t)$ von $x'$ zu $x$. Dies induziert einen Gruppenhomomorphismus 
        \begin{equation*}
            \overline{V_*}: \pi_1(X, x) \to \pi_1(X,x),
        \end{equation*}
        der invers ist zu $V_*$: Es gilt 
        \begin{align*}
            \overline{V_*}(V_*([w])) = \overline{V_*}([v * w * \overline{v}]) = [\underset{\hmtp k}{\underbrace{\overline{v} * v}} * w * \underset{\hmtp k}{\underbrace{\overline{v} * v}}] = [w].
        \end{align*}
        Analog ist $V_* \circ \overline{V_*} = \id_{\pi_1(X,x)}$, also ist $\overline{V_*}$ invers zu $V_*$.
    \end{description}
\end{proof}

\textbf{Warnung.} Der Isomorphismus 
\begin{equation*}
    V_* : \pi_1(X,x') \to \pi_1(X,x)
\end{equation*}
hängt im Allgemeinen vom Weg $v$ ab.

%%%%%%%%%%%%%%%%%%%%%%%%%%%%%%%%%%%%%%%%%%%%%%%%%%%%%%%%%%%%%%%%%%%%%%%%%%%%%%%%%%%%%%%%%%%%%%%%%%%%
%%%%%%%%%%%%%%%%%%%%%%%%%%%%%%%%%%%%%%%%%% 2024 / 05 / 28 %%%%%%%%%%%%%%%%%%%%%%%%%%%%%%%%%%%%%%%%%%
%%%%%%%%%%%%%%%%%%%%%%%%%%%%%%%%%%%%%%%%%%%%%%%%%%%%%%%%%%%%%%%%%%%%%%%%%%%%%%%%%%%%%%%%%%%%%%%%%%%%

\breakline{Beginn Vorlesung 12 (28.05.2024)}



%%%%%%%%%%%%%%%%%%%%%%%%%%%%%%%%%%%%%%%%%%%%%%%%%%%%%%%%%%%%%%%%%%%%%%%%%%%%%%%%%%%%%%%%%%%%%%%%%%%%
%%%%%%%%%%%%%%%%%%%%%%%%%%%%%%%%%%%%%%%%%% 2024 / 05 / 30 %%%%%%%%%%%%%%%%%%%%%%%%%%%%%%%%%%%%%%%%%%
%%%%%%%%%%%%%%%%%%%%%%%%%%%%%%%%%%%%%%%%%%%%%%%%%%%%%%%%%%%%%%%%%%%%%%%%%%%%%%%%%%%%%%%%%%%%%%%%%%%%

\breakline{Beginn Vorlesung 13 (30.05.2024)}



\end{document}